%\iffalse
% proof-at-the-end.dtx generated using makedtx version 1.2 (c) Nicola Talbot
% Command line args:
%   -src "proof-at-the-end\.sty=>proof-at-the-end.sty"
%   -author "Léo Colisson"
%   -doc "../doc/proof-at-the-end_doc.tex"
%   -dir "/home/leo/Documents/Informatique/Programmation/LaTeX/proof-at-the-end/"
%   proof-at-the-end
% Created on 2026/1/12 21:40
%\fi
%\iffalse
%<*package>
%% \CharacterTable
%%  {Upper-case    \A\B\C\D\E\F\G\H\I\J\K\L\M\N\O\P\Q\R\S\T\U\V\W\X\Y\Z
%%   Lower-case    \a\b\c\d\e\f\g\h\i\j\k\l\m\n\o\p\q\r\s\t\u\v\w\x\y\z
%%   Digits        \0\1\2\3\4\5\6\7\8\9
%%   Exclamation   \!     Double quote  \"     Hash (number) \#
%%   Dollar        \$     Percent       \%     Ampersand     \&
%%   Acute accent  \'     Left paren    \(     Right paren   \)
%%   Asterisk      \*     Plus          \+     Comma         \,
%%   Minus         \-     Point         \.     Solidus       \/
%%   Colon         \:     Semicolon     \;     Less than     \<
%%   Equals        \=     Greater than  \>     Question mark \?
%%   Commercial at \@     Left bracket  \[     Backslash     \\
%%   Right bracket \]     Circumflex    \^     Underscore    \_
%%   Grave accent  \`     Left brace    \{     Vertical bar  \|
%%   Right brace   \}     Tilde         \~}
%</package>
%\fi
% \iffalse
% Doc-Source file to use with LaTeX2e
% Copyright (C) 2026 Léo Colisson, all rights reserved.
% \fi
% \iffalse
%<*driver>
% Options for packages loaded elsewhere
\PassOptionsToPackage{unicode}{hyperref}
\PassOptionsToPackage{hyphens}{url}
\documentclass[
]{article}
\usepackage{xcolor}
\usepackage{amsmath,amssymb}
\setcounter{secnumdepth}{5}
\usepackage{iftex}
\ifPDFTeX
  \usepackage[T1]{fontenc}
  \usepackage[utf8]{inputenc}
  \usepackage{textcomp} % provide euro and other symbols
\else % if luatex or xetex
  \usepackage{unicode-math} % this also loads fontspec
  \defaultfontfeatures{Scale=MatchLowercase}
  \defaultfontfeatures[\rmfamily]{Ligatures=TeX,Scale=1}
\fi
\usepackage{lmodern}
\ifPDFTeX\else
  % xetex/luatex font selection
\fi
% Use upquote if available, for straight quotes in verbatim environments
\IfFileExists{upquote.sty}{\usepackage{upquote}}{}
\IfFileExists{microtype.sty}{% use microtype if available
  \usepackage[]{microtype}
  \UseMicrotypeSet[protrusion]{basicmath} % disable protrusion for tt fonts
}{}
\makeatletter
\@ifundefined{KOMAClassName}{% if non-KOMA class
  \IfFileExists{parskip.sty}{%
    \usepackage{parskip}
  }{% else
    \setlength{\parindent}{0pt}
    \setlength{\parskip}{6pt plus 2pt minus 1pt}}
}{% if KOMA class
  \KOMAoptions{parskip=half}}
\makeatother
\usepackage{color}
\usepackage{fancyvrb}
\newcommand{\VerbBar}{|}
\newcommand{\VERB}{\Verb[commandchars=\\\{\}]}
\DefineVerbatimEnvironment{Highlighting}{Verbatim}{commandchars=\\\{\}}
% Add ',fontsize=\small' for more characters per line
\newenvironment{Shaded}{}{}
\newcommand{\AlertTok}[1]{\textcolor[rgb]{1.00,0.00,0.00}{\textbf{#1}}}
\newcommand{\AnnotationTok}[1]{\textcolor[rgb]{0.38,0.63,0.69}{\textbf{\textit{#1}}}}
\newcommand{\AttributeTok}[1]{\textcolor[rgb]{0.49,0.56,0.16}{#1}}
\newcommand{\BaseNTok}[1]{\textcolor[rgb]{0.25,0.63,0.44}{#1}}
\newcommand{\BuiltInTok}[1]{\textcolor[rgb]{0.00,0.50,0.00}{#1}}
\newcommand{\CharTok}[1]{\textcolor[rgb]{0.25,0.44,0.63}{#1}}
\newcommand{\CommentTok}[1]{\textcolor[rgb]{0.38,0.63,0.69}{\textit{#1}}}
\newcommand{\CommentVarTok}[1]{\textcolor[rgb]{0.38,0.63,0.69}{\textbf{\textit{#1}}}}
\newcommand{\ConstantTok}[1]{\textcolor[rgb]{0.53,0.00,0.00}{#1}}
\newcommand{\ControlFlowTok}[1]{\textcolor[rgb]{0.00,0.44,0.13}{\textbf{#1}}}
\newcommand{\DataTypeTok}[1]{\textcolor[rgb]{0.56,0.13,0.00}{#1}}
\newcommand{\DecValTok}[1]{\textcolor[rgb]{0.25,0.63,0.44}{#1}}
\newcommand{\DocumentationTok}[1]{\textcolor[rgb]{0.73,0.13,0.13}{\textit{#1}}}
\newcommand{\ErrorTok}[1]{\textcolor[rgb]{1.00,0.00,0.00}{\textbf{#1}}}
\newcommand{\ExtensionTok}[1]{#1}
\newcommand{\FloatTok}[1]{\textcolor[rgb]{0.25,0.63,0.44}{#1}}
\newcommand{\FunctionTok}[1]{\textcolor[rgb]{0.02,0.16,0.49}{#1}}
\newcommand{\ImportTok}[1]{\textcolor[rgb]{0.00,0.50,0.00}{\textbf{#1}}}
\newcommand{\InformationTok}[1]{\textcolor[rgb]{0.38,0.63,0.69}{\textbf{\textit{#1}}}}
\newcommand{\KeywordTok}[1]{\textcolor[rgb]{0.00,0.44,0.13}{\textbf{#1}}}
\newcommand{\NormalTok}[1]{#1}
\newcommand{\OperatorTok}[1]{\textcolor[rgb]{0.40,0.40,0.40}{#1}}
\newcommand{\OtherTok}[1]{\textcolor[rgb]{0.00,0.44,0.13}{#1}}
\newcommand{\PreprocessorTok}[1]{\textcolor[rgb]{0.74,0.48,0.00}{#1}}
\newcommand{\RegionMarkerTok}[1]{#1}
\newcommand{\SpecialCharTok}[1]{\textcolor[rgb]{0.25,0.44,0.63}{#1}}
\newcommand{\SpecialStringTok}[1]{\textcolor[rgb]{0.73,0.40,0.53}{#1}}
\newcommand{\StringTok}[1]{\textcolor[rgb]{0.25,0.44,0.63}{#1}}
\newcommand{\VariableTok}[1]{\textcolor[rgb]{0.10,0.09,0.49}{#1}}
\newcommand{\VerbatimStringTok}[1]{\textcolor[rgb]{0.25,0.44,0.63}{#1}}
\newcommand{\WarningTok}[1]{\textcolor[rgb]{0.38,0.63,0.69}{\textbf{\textit{#1}}}}
\usepackage{graphicx}
\makeatletter
\newsavebox\pandoc@box
\newcommand*\pandocbounded[1]{% scales image to fit in text height/width
  \sbox\pandoc@box{#1}%
  \Gscale@div\@tempa{\textheight}{\dimexpr\ht\pandoc@box+\dp\pandoc@box\relax}%
  \Gscale@div\@tempb{\linewidth}{\wd\pandoc@box}%
  \ifdim\@tempb\p@<\@tempa\p@\let\@tempa\@tempb\fi% select the smaller of both
  \ifdim\@tempa\p@<\p@\scalebox{\@tempa}{\usebox\pandoc@box}%
  \else\usebox{\pandoc@box}%
  \fi%
}
% Set default figure placement to htbp
\def\fps@figure{htbp}
\makeatother
\setlength{\emergencystretch}{3em} % prevent overfull lines
\providecommand{\tightlist}{%
  \setlength{\itemsep}{0pt}\setlength{\parskip}{0pt}}
\usepackage{bookmark}
\IfFileExists{xurl.sty}{\usepackage{xurl}}{} % add URL line breaks if available
\urlstyle{same}
\hypersetup{
  pdftitle={Proof-at-the-end, or how to move proofs in appendix in LaTeX},
  pdfauthor={Léo Colisson},
  hidelinks,
  pdfcreator={LaTeX via pandoc}}

\title{Proof-at-the-end, or how to move proofs in appendix in LaTeX}
\author{Léo Colisson}
\date{}

\begin{document}
\DocInput{proof-at-the-end.dtx}
\end{document}
%</driver>
%\fi
%\maketitle
%
%{
%\setcounter{tocdepth}{3}
%\tableofcontents
%}
%WARNING: this package is definitely usable, but may not be as tested and
%stable as standard packages. Note also that it may change a bit in the
%future, but we will of course try to avoid as much as possible backward
%incompatibilities. To ensure your document won't change when the package
%get updates, it is always safer to copy the
%\texttt{proof-at-the-end.sty} file at the root of your project.
%
%\section{Introduction}\label{introduction}
%
%This package aims to provide a way to easily move proofs in the
%appendix. You can:
%
%\begin{itemize}
%\tightlist
%\item
%  Move proofs in different places/sections by giving different
%  ``categories'' to the theorems
%\item
%  Create links from theorem to proof, and from proof to theorem (and
%  configure the text)
%\item
%  Restate the theorem in appendix (or before)
%\item
%  Choose what you want to keep in the main text or in appendix (you can
%  put in appendix either everything, just the proof, the theorem and the
%  proof by keeping (or not) a copy in the body of the article\ldots)
%\item
%  Duplicate the proof in appendix and in the current section, practical
%  to use synctex during the proof writting
%\item
%  Add text that would appear only in the appendix (or in both body and
%  appendix)
%\item
%  Easily change the defaults, and create your own styles/environments
%\item
%  Include sketch of proof in the main text, and full proof in appendix
%\item
%  Change the text of the link, for example to translate into another
%  language
%\item
%  Have a nice environment-based commands in order to mimic the usual
%  theorem/proof structure.
%\item
%  Separate the main article and the body in two documents (some journals
%  may ask for a separate file for the appendix)
%\end{itemize}
%
%NB: This project is hosted on github at
%\url{https://github.com/leo-colisson/proof-at-the-end} . Feel free to
%contribute, report bugs, or read/copy-paste the documentation/examples
%from there.
%
%Disclaimer: This package is still in beta and not considered as stable.
%
%This package is licensed under LPPL v1.3, and the last version of this
%package on CTAN is 2019/05/21.
%
%\section{Demo}\label{demo}
%
%If you just want to see an example of what you can do, you can directly
%open the file \texttt{demo.pdf} (also available online at
%\url{https://github.com/leo-colisson/proof-at-the-end/raw/master/demo.pdf})
%to see what is possible, or generate it with
%
%\begin{Shaded}
%\begin{Highlighting}[]
%\FunctionTok{git}\NormalTok{ clone https://github.com/leo{-}colisson/proof{-}at{-}the{-}end.git}
%\ExtensionTok{pdflatex}\NormalTok{ demo.tex }\KeywordTok{\&\&} \ExtensionTok{pdflatex}\NormalTok{ demo.tex}
%\end{Highlighting}
%\end{Shaded}
%
%\section{Quickstart}\label{quickstart}
%
%\subsection{Super Quick QuickStart}\label{super-quick-quickstart}
%
%If \texttt{proof-at-the-end} it's not installed in your CTAN
%distribution, copy the
%\href{https://github.com/leo-colisson/proof-at-the-end/}{\texttt{proof-at-the-end.sty}
%file} in your project. Then, in your project, create new theorem/lemma
%environments (using any tool you like, like asmthm, ntheorem and
%thmtools), just, until
%\href{https://github.com/muzimuzhi/thmtools/issues/70}{this upstream
%\texttt{thmtool} bug} is fixed, links will be broken unless you do one
%of the following actions: - call
%\texttt{\textbackslash{}hypersetup\{hypertexnames=false\}} at the end of
%the preamble (but this may fails if you build an index), - or, possibly
%cleaner, if you use something like
%\texttt{\textbackslash{}newtheorem\{thm\}\{Theorem\}{[}section{]}}, then
%you must call
%\texttt{\textbackslash{}renewcommand\textbackslash{}theHthm\{\textbackslash{}thesection\ .\textbackslash{}arabic\ \{thm\}\}}.
%If you are unhappy with these workarounds, raise your voice in the above
%\texttt{thmtool} issue.
%
%Then, load the library using
%\texttt{\textbackslash{}usepackage{[}createShortEnv{]}\{proof-at-the-end\}}
%(note that you need the version 2022/01/28 to have
%\texttt{createShortEnv} defined, otherwise you need to use directly
%\texttt{theoremEnd}) and write your theorem using:
%
%\begin{Shaded}
%\begin{Highlighting}[]
%\KeywordTok{\textbackslash{}begin}\NormalTok{\{}\ExtensionTok{thmE}\NormalTok{\}[My title][end, restate]}
%\NormalTok{  I am a restated theorem whose proof goes in appendix (compile me twice).}
%\KeywordTok{\textbackslash{}end}\NormalTok{\{}\ExtensionTok{thmE}\NormalTok{\}}
%\KeywordTok{\textbackslash{}begin}\NormalTok{\{}\ExtensionTok{proofE}\NormalTok{\}}
%\NormalTok{  And I am a proof.}
%\KeywordTok{\textbackslash{}end}\NormalTok{\{}\ExtensionTok{proofE}\NormalTok{\}}
%\end{Highlighting}
%\end{Shaded}
%
%The options (here \texttt{restate} and \texttt{end}) tells what should
%go in appendix, how to configure the links\ldots{} Here the proof should
%go in appendix, and the \texttt{restate} options states that the theorem
%should be restated before the proof. Note that in case you don't want
%any title but do want to add options, make sure to insert an empty pair
%of brackets like in
%\texttt{\textbackslash{}begin\{thmE\}{[}{]}{[}end,\ text\ link={]}} (the
%option \texttt{text\ link} is used to choose the text of the link to the
%proof, writting \texttt{text\ link=} is used to display no link at all).
%
%You can find below a full example to compile which should produce this
%output (sorry, this is a screenshot, that way I can export it easily in
%the github page). Just make sure to compile twice to get references.
%
%\begin{figure}
%\centering
%\pandocbounded{\includegraphics[keepaspectratio,alt={Screenshot of the quickstart}]{screenshot.png}}
%\caption{Screenshot of the quickstart}
%\end{figure}
%
%You can directly copy them (with a more complete demo) in the github
%page here https://github.com/leo-colisson/proof-at-the-end/.
%
%\begin{Shaded}
%\begin{Highlighting}[]
%\BuiltInTok{\textbackslash{}documentclass}\NormalTok{\{}\ExtensionTok{article}\NormalTok{\}}
%
%\BuiltInTok{\textbackslash{}usepackage}\NormalTok{\{}\ExtensionTok{xparse}\NormalTok{\} }\CommentTok{\% Not needed with recent LaTeX}
%\BuiltInTok{\textbackslash{}usepackage}\NormalTok{\{}\ExtensionTok{amsthm}\NormalTok{\}}
%
%\CommentTok{\% Create new theorems, or use ntheorem/thmtools/...}
%\FunctionTok{\textbackslash{}newtheorem}\NormalTok{\{thm\}\{Theorem\}[section]}
%\CommentTok{\% Default uses autoref, but you can also use cleveref, see the documentation.}
%\CommentTok{\% Say to autoref that "thm" are Theorems.}
%\FunctionTok{\textbackslash{}newcommand}\NormalTok{\{}\ExtensionTok{\textbackslash{}thmautorefname}\NormalTok{\}\{Theorem\}}
%
%\CommentTok{\%\% Load the library. createShortEnv automatically creates the shortcuts}
%\CommentTok{\%\% thmE, theoremE, lemmaE, propositionE, corollaryE, proofE. See \textbackslash{}newEndThm for more details.}
%\BuiltInTok{\textbackslash{}usepackage}\NormalTok{[createShortEnv]\{}\ExtensionTok{proof{-}at{-}the{-}end}\NormalTok{\}}
%
%\KeywordTok{\textbackslash{}begin}\NormalTok{\{}\ExtensionTok{document}\NormalTok{\}}
%
%\KeywordTok{\textbackslash{}section}\NormalTok{\{Theorems\}}
%
%\CommentTok{\%\% If you don\textquotesingle{}t have any title, leave an empty bracket like \textbackslash{}begin\{thmE\}[][end]}
%\KeywordTok{\textbackslash{}begin}\NormalTok{\{}\ExtensionTok{thmE}\NormalTok{\}[My title][end, restate]}
%\NormalTok{  I am a theorem}
%\KeywordTok{\textbackslash{}end}\NormalTok{\{}\ExtensionTok{thmE}\NormalTok{\}}
%\KeywordTok{\textbackslash{}begin}\NormalTok{\{}\ExtensionTok{proofE}\NormalTok{\}}
%\NormalTok{  And I am a proof.}
%\KeywordTok{\textbackslash{}end}\NormalTok{\{}\ExtensionTok{proofE}\NormalTok{\}}
%
%\KeywordTok{\textbackslash{}section}\NormalTok{\{Proofs\}}
%\FunctionTok{\textbackslash{}printProofs}
%
%\KeywordTok{\textbackslash{}end}\NormalTok{\{}\ExtensionTok{document}\NormalTok{\}}
%\end{Highlighting}
%\end{Shaded}
%
%\subsection{Install}\label{install}
%
%If your CTAN distribution is recent enough, you have nothing to do (if
%you are using overleaf, note that you can configure your project to use
%a more recent TexLive distribution: this packages is included starting
%from TexLive 2019). Otherwise if it's not yet in your CTAN distribution,
%first download the \texttt{proof-at-the-end.sty} file and insert it in
%the root of your project with the following commands on unix (you can
%also clone this repository if you prefer, or just manually download or
%copy/paste the files on Windows).
%
%\begin{Shaded}
%\begin{Highlighting}[]
%\BuiltInTok{cd} \OperatorTok{\textless{}}\NormalTok{your project}\OperatorTok{\textgreater{}}
%\VariableTok{repopratend}\OperatorTok{=}\StringTok{"https://raw.githubusercontent.com/leo{-}colisson/proof{-}at{-}the{-}end"}
%\FunctionTok{wget} \VariableTok{$\{repopratend\}}\NormalTok{/master/proof{-}at{-}the{-}end.sty}
%\end{Highlighting}
%\end{Shaded}
%
%It also requires a recent version of xparse (not present in TexLive 2018
%and before basically), otherwise you will get an error like:
%
%\begin{verbatim}
%Unknown argument type 'b' replaced by 'm'
%\end{verbatim}
%
%While
%\href{https://raw.githubusercontent.com/leo-colisson/proof-at-the-end/c2150f4feba4709d47e43c00e5f8eb93ce345cab/xparse.sty}{we
%used to provide} for simplicity a \texttt{xparse.sty} file (version
%\texttt{2018-04-12}) in older versions of this repository, we removed it
%to
%\href{https://github.com/leo-colisson/proof-at-the-end/issues/6}{avoid
%license issues and conflicts with recent LaTeX distributions that
%already include xparse}. Nevertheless, if you are using an older
%distribution, you can also download \texttt{xparse} manually
%\href{http://mirrors.ctan.org/install/macros/latex/contrib/l3packages.tds.zip}{here},
%unzip it, and copy the \texttt{.sty} files contained in
%\texttt{tex/latex/l3packages/xparse/} at the root of your project. On
%linux, you can do that via:
%
%\begin{Shaded}
%\begin{Highlighting}[]
%\BuiltInTok{cd} \OperatorTok{\textless{}}\NormalTok{your project}\OperatorTok{\textgreater{}}
%\FunctionTok{wget}\NormalTok{ http://mirrors.ctan.org/install/macros/latex/contrib/l3packages.tds.zip}
%\FunctionTok{unzip} \AttributeTok{{-}d}\NormalTok{ . }\AttributeTok{{-}j}\NormalTok{ l3packages.tds.zip }\StringTok{\textquotesingle{}tex/latex/l3packages/xparse/*\textquotesingle{}}
%\FunctionTok{rm} \AttributeTok{{-}f}\NormalTok{ l3packages.tds.zip xparse.ltx}
%\end{Highlighting}
%\end{Shaded}
%
%If you have an old distribution of LaTeX (2017 or before basically), you
%may also need a
%\href{https://tex.stackexchange.com/questions/489646/expl3-and-recent-xparse-on-overleaf-no-expl3-loader-detected/489649?noredirect=1\#comment1236409_489649}{more
%recent expl3} or you will get an error like:
%
%\begin{verbatim}
%Package xparse Error: Support package l3kernel too old.
%\end{verbatim}
%
%To install it, download the zip files
%\url{http://mirrors.ctan.org/install/macros/latex/contrib/l3kernel.tds.zip}
%and
%\url{http://mirrors.ctan.org/install/macros/latex/contrib/l3backend.tds.zip}
%(\texttt{l3backend} is
%\href{https://tex.stackexchange.com/questions/499082/file-l3backend-pdfmode-def-not-found-when-loading-expl3}{since
%July 1st, 2019} a dependency of l3kernel), unzip, and copy the content
%of the directories \texttt{tex/latex/l3kernel/} (first archive) and
%\texttt{tex/latex/l3backend/} (second archive) into your project. On
%linux it's a matter of two commands in your project:
%
%\begin{Shaded}
%\begin{Highlighting}[]
%\BuiltInTok{cd} \OperatorTok{\textless{}}\NormalTok{your project}\OperatorTok{\textgreater{}}
%\FunctionTok{wget}\NormalTok{ http://mirrors.ctan.org/install/macros/latex/contrib/l3kernel.tds.zip}
%\FunctionTok{unzip} \AttributeTok{{-}d}\NormalTok{ . }\AttributeTok{{-}j}\NormalTok{ l3kernel.tds.zip }\StringTok{\textquotesingle{}tex/latex/l3kernel/*\textquotesingle{}}
%\FunctionTok{rm}\NormalTok{ l3kernel.tds.zip}
%\FunctionTok{wget}\NormalTok{ http://mirrors.ctan.org/install/macros/latex/contrib/l3backend.tds.zip}
%\FunctionTok{unzip} \AttributeTok{{-}d}\NormalTok{ . }\AttributeTok{{-}j}\NormalTok{ l3backend.tds.zip }\StringTok{\textquotesingle{}tex/latex/l3backend/*\textquotesingle{}}
%\FunctionTok{rm}\NormalTok{ l3backend.tds.zip}
%\end{Highlighting}
%\end{Shaded}
%
%If you don't want to pollute your main project with all these files, you
%may be interested to put them in a subfolder and update the environment
%variable \texttt{TEXINPUTS} or, if you use latexmk or overleaf, you can
%write instead a \texttt{latexmkrc} file as explained
%\href{https://www.overleaf.com/learn/latex/Questions/I_have_a_lot_of_.cls,_.sty,_.bst_files,_and_I_want_to_put_them_in_a_folder_to_keep_my_project_uncluttered._But_my_project_is_not_finding_them_to_compile_correctly}{here}.
%
%If you still have errors, it's likely that your LaTeX distribution is
%just too old (LaTeX 2016 seems to be too old to install a recent l3exp).
%So just upgrade to a decent version, and you should not have any
%troubles anymore.
%
%You can test your install by compiling the \texttt{demo.tex} file
%present in the repository. When it compiles, you are ready to start!
%
%\subsection{Use in your project}\label{use-in-your-project}
%
%Now, you can load the library in your project by simply using:
%
%\begin{Shaded}
%\begin{Highlighting}[]
%\BuiltInTok{\textbackslash{}usepackage}\NormalTok{\{}\ExtensionTok{proof{-}at{-}the{-}end}\NormalTok{\}}
%\end{Highlighting}
%\end{Shaded}
%
%Then, you can configure your theorem/lemma/\ldots{} environments as
%usual, by using any counter you like\ldots:
%
%\begin{Shaded}
%\begin{Highlighting}[]
%\BuiltInTok{\textbackslash{}usepackage}\NormalTok{\{}\ExtensionTok{amssymb, amsthm, amsmath}\NormalTok{\}}
%\CommentTok{\% Theorems}
%\FunctionTok{\textbackslash{}newtheorem}\NormalTok{\{thm\}\{Theorem\}[section]}
%\FunctionTok{\textbackslash{}newtheorem*}\NormalTok{\{thm*\}\{Theorem\}}
%\FunctionTok{\textbackslash{}providecommand*\textbackslash{}thmautorefname}\NormalTok{\{Theorem\}}
%\CommentTok{\% Lemmata}
%\FunctionTok{\textbackslash{}newtheorem}\NormalTok{\{lemma\}[thm]\{Lemma\}}
%\FunctionTok{\textbackslash{}newtheorem*}\NormalTok{\{lemma*\}\{Lemma\}}
%\FunctionTok{\textbackslash{}providecommand*\textbackslash{}lemmaautorefname}\NormalTok{\{Lemma\}}
%\end{Highlighting}
%\end{Shaded}
%
%And inside your document, you can use the following syntax to create a
%new theorem:
%
%\begin{Shaded}
%\begin{Highlighting}[]
%\KeywordTok{\textbackslash{}begin}\NormalTok{\{}\ExtensionTok{theoremEnd}\NormalTok{\}[OPTIONS]\{THEOREM ENVIRONMENT\}[OPTIONAL TITLE]}
%\NormalTok{    YOUR THEOREM, with eventually labels like }\KeywordTok{\textbackslash{}label}\NormalTok{\{}\ExtensionTok{thm:OPTIONAL LABEL}\NormalTok{\}}
%\KeywordTok{\textbackslash{}end}\NormalTok{\{}\ExtensionTok{theoremEnd}\NormalTok{\}}
%\KeywordTok{\textbackslash{}begin}\NormalTok{\{}\ExtensionTok{proofEnd}\NormalTok{\} }\CommentTok{\%\% Optional environment}
%\NormalTok{    YOUR (OPTIONAL) PROOF}
%\KeywordTok{\textbackslash{}end}\NormalTok{\{}\ExtensionTok{proofEnd}\NormalTok{\}}
%\end{Highlighting}
%\end{Shaded}
%
%For example:
%
%\begin{Shaded}
%\begin{Highlighting}[]
%\KeywordTok{\textbackslash{}begin}\NormalTok{\{}\ExtensionTok{theoremEnd}\NormalTok{\}\{thm\}[Yes I can have a title]}
%  \KeywordTok{\textbackslash{}label}\NormalTok{\{}\ExtensionTok{thm:ilikelabels}\NormalTok{\}}
%\NormalTok{  Creating a new theorem is easy}
%\KeywordTok{\textbackslash{}end}\NormalTok{\{}\ExtensionTok{theoremEnd}\NormalTok{\}}
%\KeywordTok{\textbackslash{}begin}\NormalTok{\{}\ExtensionTok{proofEnd}\NormalTok{\}}
%\NormalTok{  You want a proof? Here is it!}
%\KeywordTok{\textbackslash{}end}\NormalTok{\{}\ExtensionTok{proofEnd}\NormalTok{\}}
%\end{Highlighting}
%\end{Shaded}
%
%Since \texttt{theoremEnd} may be a bit verbose to use, we provide since
%the version 2022/01/28, an option: if you load the program using
%\texttt{\textbackslash{}usepackage{[}createShortEnv{]}\{proof-at-the-end\}},
%you will automatically have the shortcut environments \texttt{proofE}
%(to replace \texttt{proof}), \texttt{thmE} (to replace \texttt{thm}),
%\texttt{theoremE} (to replace\ldots{} ok you got the pattern),
%\texttt{lemmaE}, \texttt{propositionE} and \texttt{corollaryE}. Note
%that it is your role to define the initial commands \texttt{thm}\ldots,
%but you don't need to define them all, only those that you use. Then,
%you can use them like:
%
%\begin{Shaded}
%\begin{Highlighting}[]
%\KeywordTok{\textbackslash{}begin}\NormalTok{\{}\ExtensionTok{thmE}\NormalTok{\}[My title][end, restate]}
%\NormalTok{  I am a restated theorem whose proof goes in appendix (compile me twice).}
%\KeywordTok{\textbackslash{}end}\NormalTok{\{}\ExtensionTok{thmE}\NormalTok{\}}
%\KeywordTok{\textbackslash{}begin}\NormalTok{\{}\ExtensionTok{proofE}\NormalTok{\}}
%\NormalTok{  And I am a proof.}
%\KeywordTok{\textbackslash{}end}\NormalTok{\{}\ExtensionTok{proofE}\NormalTok{\}}
%\end{Highlighting}
%\end{Shaded}
%
%Finally, to display the proofs, just use
%\texttt{\textbackslash{}printProofs} where you like:
%
%\begin{Shaded}
%\begin{Highlighting}[]
%\FunctionTok{\textbackslash{}printProofs}
%\end{Highlighting}
%\end{Shaded}
%
%If you would like to display a lemma instead, just change
%\texttt{\{thm\}} into \texttt{\{lemma\}}, or into any other theorem
%environment you defined! You can now compile safely your document ;)
%
%NB: if you want to make sure all the references are linked correctly,
%make sure to compile twice the document!
%
%Isn't it simple ?
%
%\subsection{\texorpdfstring{A note on
%\texttt{hyperref}}{A note on hyperref}}\label{a-note-on-hyperref}
%
%An older version of the package explicitely loaded \texttt{hyperref}
%with no option, and an option clash could occur if the user wanted to
%add options to \texttt{hyperref}. We now push the loading to the very
%end to avoid that issue, but if you want to load packages \emph{after}
%\texttt{hyperref} (like \texttt{cleveref}), it is your job to include
%\texttt{hyperref}, for example by putting at the end of the preambule
%something like:
%
%\begin{Shaded}
%\begin{Highlighting}[]
%\BuiltInTok{\textbackslash{}usepackage}\NormalTok{[colorlinks]\{}\ExtensionTok{hyperref}\NormalTok{\}}
%\BuiltInTok{\textbackslash{}usepackage}\NormalTok{\{}\ExtensionTok{cleveref}\NormalTok{\}}
%\end{Highlighting}
%\end{Shaded}
%
%\section{Use cases}\label{use-cases}
%
%\subsection{Configuration and how to use and create
%styles}\label{configuration-and-how-to-use-and-create-styles}
%
%You can very easily configure this package, and choose how each
%theorem/proof must be displayed by providing a value in
%\texttt{OPTIONS}. For example, if you would like to keep the proof of a
%theorem in the main text like any normal theorem, use the
%\texttt{normal} option.
%
%\begin{Shaded}
%\begin{Highlighting}[]
%\KeywordTok{\textbackslash{}begin}\NormalTok{\{}\ExtensionTok{theoremEnd}\NormalTok{\}[normal]\{thm\}[A title]}
%\NormalTok{  You can easily turn a theorem back into a normal theorem!}
%\KeywordTok{\textbackslash{}end}\NormalTok{\{}\ExtensionTok{theoremEnd}\NormalTok{\}}
%\KeywordTok{\textbackslash{}begin}\NormalTok{\{}\ExtensionTok{proofEnd}\NormalTok{\}}
%\NormalTok{  And keep the proof with you!}
%\KeywordTok{\textbackslash{}end}\NormalTok{\{}\ExtensionTok{proofEnd}\NormalTok{\}}
%\end{Highlighting}
%\end{Shaded}
%
%If you loaded the package using
%\texttt{\textbackslash{}usepackage{[}createShortEnv{]}\{proof-at-the-end\}}
%(starting from 2022/01/28), you also have preconfigured shortcuts for
%multiple environments, namely \texttt{proofE} (to replace
%\texttt{proof}), \texttt{thmE} (to replace \texttt{thm}),
%\texttt{theoremE} (to replace\ldots{} ok you got the pattern),
%\texttt{lemmaE}, \texttt{propositionE}, \texttt{propertyE},
%\texttt{factE} and \texttt{corollaryE}. This way, you can simplify the
%above as:
%
%\begin{Shaded}
%\begin{Highlighting}[]
%\KeywordTok{\textbackslash{}begin}\NormalTok{\{}\ExtensionTok{thmE}\NormalTok{\}[A title][normal]}
%\NormalTok{  You can easily turn a theorem back into a normal theorem!}
%\KeywordTok{\textbackslash{}end}\NormalTok{\{}\ExtensionTok{thmE}\NormalTok{\}}
%\KeywordTok{\textbackslash{}begin}\NormalTok{\{}\ExtensionTok{proofE}\NormalTok{\}}
%\NormalTok{  And keep the proof with you!}
%\KeywordTok{\textbackslash{}end}\NormalTok{\{}\ExtensionTok{proofE}\NormalTok{\}}
%\end{Highlighting}
%\end{Shaded}
%
%(You can of course remove the \texttt{{[}...{]}} containing the options
%and the title, however if you want to add options without any title,
%\textbf{make sure to use empty brackets for the title}, like
%\texttt{\textbackslash{}begin\{thmE\}{[}{]}{[}normal{]}}.
%
%Since 2022/01/28, you can also create new theorem environment that
%directly come with a preconfigured style (you may also find
%\texttt{\textbackslash{}pratendSetLocal} useful if you want to change
%the style in a whole section). For instance to create a
%\texttt{thmNormal} environment having the \texttt{normal} option and a
%\texttt{thmSpecial} environment having the
%\texttt{all\ end,\ category=mySpecial} options, just do:
%
%\begin{verbatim}
%\newEndThm[normal]{thmNormal}{thm} % 'thm' is the name of the existing thm environment 
%\newEndThm[all end, category=mySpecial]{thmSpecial}{thm}
%\end{verbatim}
%
%This way, you can simply use
%\texttt{\textbackslash{}begin\{thmSpecial\}{[}Special\ title{]}\ ...\ \textbackslash{}end\{thmSpecial\}}
%to create a theorem going in the special category (see below). (you also
%have \texttt{\textbackslash{}newEndProof{[}options{]}\{newProofEnv\}} to
%define new proofs, but it is less useful as options are typically given
%in the theorem).
%
%The options are in fact a set of keys/values, thanks to
%\texttt{pgfkeys}. So you can combine them with comma separated list like
%that (order matters, as the right-most values may overwrite
%configuration set by left-most values):
%
%\begin{Shaded}
%\begin{Highlighting}[]
%\KeywordTok{\textbackslash{}begin}\NormalTok{\{}\ExtensionTok{theoremEnd}\NormalTok{\}[proof at the end,}
%\NormalTok{                   no link to proof,}
%\NormalTok{                   text proof=\{Difficult proof\}}
%\NormalTok{                  ]\{thm\}[A title]}
%\NormalTok{  Each theorem can have a custom configuration!}
%\KeywordTok{\textbackslash{}end}\NormalTok{\{}\ExtensionTok{theoremEnd}\NormalTok{\}}
%\KeywordTok{\textbackslash{}begin}\NormalTok{\{}\ExtensionTok{proofEnd}\NormalTok{\}}
%\NormalTok{  Quite practical, isn\textquotesingle{}t it?}
%\KeywordTok{\textbackslash{}end}\NormalTok{\{}\ExtensionTok{proofEnd}\NormalTok{\}}
%\end{Highlighting}
%\end{Shaded}
%
%You can easily create your own styles like that:
%
%\begin{Shaded}
%\begin{Highlighting}[]
%\FunctionTok{\textbackslash{}pgfkeys}\NormalTok{\{/prAtEnd/my great style/.style=\{}
%\NormalTok{    proot at the end,}
%\NormalTok{    no link to proof,}
%\NormalTok{    text proof=\{Difficult proof\},}
%\NormalTok{  \}}
%\NormalTok{\}}
%\end{Highlighting}
%\end{Shaded}
%
%You can also change the default configuration when you load the package
%by nesting the configuration into a \texttt{conf} key:
%
%\begin{Shaded}
%\begin{Highlighting}[]
%\BuiltInTok{\textbackslash{}usepackage}\NormalTok{[conf=\{normal, one big link\}]\{}\ExtensionTok{proof{-}at{-}the{-}end}\NormalTok{\}}
%\end{Highlighting}
%\end{Shaded}
%
%Note however that for now it is \emph{not} possible to use macros
%directly inside the options when you load the package, so if you need to
%use more complicated configuration, you can overwrite the
%\texttt{global\ custom\ defaults} style for global configuration, and
%the \texttt{local\ custom\ defaults} style for local configuration
%(useful for example if you want to define a category for a single
%section):
%
%\begin{Shaded}
%\begin{Highlighting}[]
%\FunctionTok{\textbackslash{}pgfkeys}\NormalTok{\{/prAtEnd/global custom defaults/.style=\{}
%\NormalTok{    one big link=\{Go to proof on page\textasciitilde{}}\FunctionTok{\textbackslash{}pageref*}\NormalTok{\{proof:prAtEnd}\FunctionTok{\textbackslash{}pratendcountercurrent}\NormalTok{\}\}}
%\NormalTok{  \}}
%\NormalTok{\}}
%\end{Highlighting}
%\end{Shaded}
%
%or, for short (needs version \textgreater{} 2022/02/04):
%
%\begin{Shaded}
%\begin{Highlighting}[]
%\FunctionTok{\textbackslash{}pratendSetGlobal}\NormalTok{\{one big link="blabla"\}}
%\end{Highlighting}
%\end{Shaded}
%
%and for local configuration:
%
%\begin{Shaded}
%\begin{Highlighting}[]
%\FunctionTok{\textbackslash{}pgfkeys}\NormalTok{\{/prAtEnd/local custom defaults/.style=\{}
%\NormalTok{    category=greattheorem}
%\NormalTok{  \}}
%\NormalTok{\}}
%\end{Highlighting}
%\end{Shaded}
%
%for short (needs version \textgreater{} 2022/02/04):
%
%\begin{Shaded}
%\begin{Highlighting}[]
%\FunctionTok{\textbackslash{}pratendSetLocal}\NormalTok{\{category=greattheorem\}}
%\end{Highlighting}
%\end{Shaded}
%
%Then you can use like that:
%
%\begin{Shaded}
%\begin{Highlighting}[]
%\KeywordTok{\textbackslash{}begin}\NormalTok{\{}\ExtensionTok{thmE}\NormalTok{\}[Title]}
%\NormalTok{  Here is a normal theorem with the proof in the main text.}
%\KeywordTok{\textbackslash{}end}\NormalTok{\{}\ExtensionTok{thmE}\NormalTok{\}}
%\KeywordTok{\textbackslash{}begin}\NormalTok{\{}\ExtensionTok{proofE}\NormalTok{\}}
%\NormalTok{  The (optional) proof}
%\KeywordTok{\textbackslash{}end}\NormalTok{\{}\ExtensionTok{proofE}\NormalTok{\}}
%  
%\KeywordTok{\textbackslash{}begin}\NormalTok{\{}\ExtensionTok{thmE}\NormalTok{\}[Title][end]}
%\NormalTok{  Here is a theorem whose proof goes to the end.}
%\KeywordTok{\textbackslash{}end}\NormalTok{\{}\ExtensionTok{thmE}\NormalTok{\}}
%\KeywordTok{\textbackslash{}begin}\NormalTok{\{}\ExtensionTok{proofE}\NormalTok{\}}
%\NormalTok{  The proof}
%\KeywordTok{\textbackslash{}end}\NormalTok{\{}\ExtensionTok{proofE}\NormalTok{\}}
%
%\KeywordTok{\textbackslash{}begin}\NormalTok{\{}\ExtensionTok{thmE}\NormalTok{\}[Title][all end]}
%\NormalTok{  Here is a theorem that goes with the proof at the end.}
%\KeywordTok{\textbackslash{}end}\NormalTok{\{}\ExtensionTok{thmE}\NormalTok{\}}
%\KeywordTok{\textbackslash{}begin}\NormalTok{\{}\ExtensionTok{proofE}\NormalTok{\}}
%\NormalTok{  The proof}
%\KeywordTok{\textbackslash{}end}\NormalTok{\{}\ExtensionTok{proofE}\NormalTok{\}}
%\end{Highlighting}
%\end{Shaded}
%
%Note also that it is possible to give options to the \texttt{proofEnd}
%environment, but it is usually useless, as it will automatically pick
%the parameters from the last \texttt{theoremEnd} environment. However,
%if for some reasons you want to change the options of the proof only,
%you can do it, but do it as your own risks ;)
%
%\subsection{Usual styles}\label{usual-styles}
%
%We predefined some pretty common styles/options. The full list is at the
%end of this document, but here is a list of the most practical ones:
%
%\begin{itemize}
%\tightlist
%\item
%  \texttt{normal}: turn the theorem into a ``normal'' theorem, with the
%  proof in the main text and nothing in appendix
%\item
%  \texttt{category=yourowncategory}: change the category of the theorem
%  (see next sub-section)
%\item
%  \texttt{end}: put the proof in appendix
%\item
%  \texttt{all\ end}: put both the theorem and the proof in appendix
%\item
%  \texttt{debug}: make sure the proof is written in the main text as
%  well. Practical when you write the proof to be able to use synctex (if
%  you use synctex with the proof in appendix, your will be unfortunately
%  moved to a temporary file that this library is using\ldots{} so
%  \textbf{make sure you don't modify the files named like
%  \texttt{*prattheenddefaultcategory.tex} or all your changes will be
%  lost at the next compilation}!).
%\item
%  \texttt{one\ big\ link}: if you prefer to have a single big link
%  instead of two links (one for the proof, one for the page)
%\item
%  \texttt{one\ big\ link\ translated=Your\ translation}: to
%  change/translate the text of the link easily
%\item
%  \texttt{text\ link\ section}: put a link looking like ``See proof in
%  section XX.''. You can use \texttt{text\ link=} to remove this link.
%\item
%  \texttt{text\ link\ section\ full\ proof}: put a link looking like
%  ``See full proof in section XX.''
%\item
%  \texttt{text\ proof\ translated=Your\ translation}: to
%  change/translate the text of the proof at the end easily
%\item
%  \texttt{global\ custom\ defaults}: empty style that you can modify to
%  change the configuration (globally)
%\item
%  \texttt{local\ custom\ defaults}: empty style that you can modify to
%  change the configuration (locally). Practical to set a category for a
%  single section.
%\end{itemize}
%
%\subsection{Categories, or how to move proofs in different
%sections}\label{categories-or-how-to-move-proofs-in-different-sections}
%
%Let's imagine that you have some proofs that are easy to do, and some
%proofs that are long but interesting. You may want to put the easy
%proofs in a different place that the long proofs. It is super easy to
%do, you just need to give a category name to the option
%\texttt{category} like here:
%
%\begin{Shaded}
%\begin{Highlighting}[]
%\KeywordTok{\textbackslash{}begin}\NormalTok{\{}\ExtensionTok{theoremEnd}\NormalTok{\}[category=mylongproofs]\{thm\}[A title]}
%\NormalTok{  You can easily change the place of the proofs}
%\KeywordTok{\textbackslash{}end}\NormalTok{\{}\ExtensionTok{theoremEnd}\NormalTok{\}}
%\KeywordTok{\textbackslash{}begin}\NormalTok{\{}\ExtensionTok{proofEnd}\NormalTok{\}}
%\NormalTok{  Just use a different category name!}
%\KeywordTok{\textbackslash{}end}\NormalTok{\{}\ExtensionTok{proofEnd}\NormalTok{\}}
%\end{Highlighting}
%\end{Shaded}
%
%and give this category name to \texttt{\textbackslash{}printProofs} in
%the section where you would like to display the proofs:
%
%\begin{Shaded}
%\begin{Highlighting}[]
%\FunctionTok{\textbackslash{}printProofs}\NormalTok{[mylongproofs]}
%\end{Highlighting}
%\end{Shaded}
%
%\subsection{Comments}\label{comments}
%
%You can also move some text in the appendix by using:
%
%\begin{Shaded}
%\begin{Highlighting}[]
%\FunctionTok{\textbackslash{}textEnd}\NormalTok{\{Your text that should go in appendix\}}
%\end{Highlighting}
%\end{Shaded}
%
%You can also give it a category as explained above, or configure it to
%be displayed in both the main text and at the end of the file with:
%
%\begin{Shaded}
%\begin{Highlighting}[]
%\FunctionTok{\textbackslash{}textEnd}\NormalTok{[both]\{I am a comment that is written in both the main text}
%\NormalTok{and the appendix\}}
%\end{Highlighting}
%\end{Shaded}
%
%You can also use the environment notation like that:
%
%\begin{Shaded}
%\begin{Highlighting}[]
%\KeywordTok{\textbackslash{}begin}\NormalTok{\{}\ExtensionTok{textAtEnd}\NormalTok{\}[options]}
%\NormalTok{  You can also use the environment syntax.}
%\KeywordTok{\textbackslash{}end}\NormalTok{\{}\ExtensionTok{textAtEnd}\NormalTok{\}}
%\end{Highlighting}
%\end{Shaded}
%
%\subsection{Restate a theorem}\label{restate-a-theorem}
%
%It is easy to restate a theorem in the appendix, to have both the
%theorem in the main text and in the appendix: just use the option
%\texttt{restate}:
%
%\begin{Shaded}
%\begin{Highlighting}[]
%\KeywordTok{\textbackslash{}begin}\NormalTok{\{}\ExtensionTok{theoremEnd}\NormalTok{\}[end, restate]\{thm\}[A title]}
%\NormalTok{  This theorem will be displayed both in main text and appendix.}
%\KeywordTok{\textbackslash{}end}\NormalTok{\{}\ExtensionTok{theoremEnd}\NormalTok{\}}
%\KeywordTok{\textbackslash{}begin}\NormalTok{\{}\ExtensionTok{proofEnd}\NormalTok{\}}
%\NormalTok{  Just use restate option.}
%\KeywordTok{\textbackslash{}end}\NormalTok{\{}\ExtensionTok{proofEnd}\NormalTok{\}}
%\end{Highlighting}
%\end{Shaded}
%
%You can also use the option \texttt{restate\ command=yourcustomcommand}
%in order to create a macro \texttt{\textbackslash{}yourcustomcommand}
%that will restate the theorem wherever you want (but after the
%definition).
%
%If you want to (re)state a theorem \emph{before} its definition (say in
%the introduction), there is also a special environment
%\texttt{theoremEndRestateBefore} that requires a (unique) custom name
%that you need to provide also later on in place of the real theorem with
%the option \texttt{restated\ before}:
%
%\begin{Shaded}
%\begin{Highlighting}[]
%\KeywordTok{\textbackslash{}section}\NormalTok{\{Introduction\}}
%\KeywordTok{\textbackslash{}begin}\NormalTok{\{}\ExtensionTok{theoremEndRestateBefore}\NormalTok{\}\{thm\}[Title]\{anamethatisusedtorestate\}}
%\NormalTok{  It is possible to state the theorem before}
%\NormalTok{  in the introduction, and restate it later}
%\KeywordTok{\textbackslash{}end}\NormalTok{\{}\ExtensionTok{theoremEndRestateBefore}\NormalTok{\}}
%
%\KeywordTok{\textbackslash{}section}\NormalTok{\{Real definition\}}
%\KeywordTok{\textbackslash{}begin}\NormalTok{\{}\ExtensionTok{theoremEnd}\NormalTok{\}[restated before]\{thm\}}
%\NormalTok{  anamethatisusedtorestate}
%\KeywordTok{\textbackslash{}end}\NormalTok{\{}\ExtensionTok{theoremEnd}\NormalTok{\}}
%\KeywordTok{\textbackslash{}begin}\NormalTok{\{}\ExtensionTok{proofEnd}\NormalTok{\}}
%\NormalTok{  Proof of the theorem, put in place of the theorem the unique name}
%\KeywordTok{\textbackslash{}end}\NormalTok{\{}\ExtensionTok{proofEnd}\NormalTok{\}}
%\end{Highlighting}
%\end{Shaded}
%
%\subsection{Translate the links}\label{translate-the-links}
%
%The more powerful way to change the text of the links is to redefine
%\texttt{text\ link} and \texttt{text\ proof} (see section List of
%options for more details). However we defined also some easy way to
%redefine the text using \texttt{one\ big\ link\ translated} and
%\texttt{text\ proof\ translated}. For example, to create your
%\texttt{french} style you can do:
%
%\begin{Shaded}
%\begin{Highlighting}[]
%\FunctionTok{\textbackslash{}pgfkeys}\NormalTok{\{/prAtEnd/french/.style=\{}
%\NormalTok{    one big link translated=\{Voir preuve page\},}
%\NormalTok{    text proof translated=\{Preuve du\}}
%\NormalTok{  \}}
%\NormalTok{\}}
%\end{Highlighting}
%\end{Shaded}
%
%\subsection{Write a sketch of proof in the main
%text}\label{write-a-sketch-of-proof-in-the-main-text}
%
%You can include a sketch of proof in the main text by simply adding a
%proof in between \texttt{theoremEnd} and \texttt{proofEnd}. An alias
%option \texttt{see\ full\ proof} can also be used to change the link
%into ``See full proof on page X.'':
%
%\begin{Shaded}
%\begin{Highlighting}[]
%\KeywordTok{\textbackslash{}begin}\NormalTok{\{}\ExtensionTok{theoremEnd}\NormalTok{\}[see full proof]\{thm\}}
%\NormalTok{  I can also write a sketch of proof, and put the full proof in appendix.}
%\KeywordTok{\textbackslash{}end}\NormalTok{\{}\ExtensionTok{theoremEnd}\NormalTok{\}}
%\KeywordTok{\textbackslash{}begin}\NormalTok{\{}\ExtensionTok{proof}\NormalTok{\}}
%\NormalTok{  Hint: look at the alias options.}
%\KeywordTok{\textbackslash{}end}\NormalTok{\{}\ExtensionTok{proof}\NormalTok{\}}
%\KeywordTok{\textbackslash{}begin}\NormalTok{\{}\ExtensionTok{proofEnd}\NormalTok{\}}
%\NormalTok{  You just use \textasciigrave{}\textasciigrave{}see full proof\textquotesingle{}\textquotesingle{} as an option}
%\KeywordTok{\textbackslash{}end}\NormalTok{\{}\ExtensionTok{proofEnd}\NormalTok{\}}
%\end{Highlighting}
%\end{Shaded}
%
%\subsection{Split the article in two documents, one for the main body
%and one for the
%appendix}\label{split-the-article-in-two-documents-one-for-the-main-body-and-one-for-the-appendix}
%
%Since 2022/02/04, it is possible to include the appendix in a separate
%file. Internally, this library creates a new file having the form
%\texttt{NAME\_OF\_FILE-pratendNAME\_CATEGORY.tex} containing the proofs
%to include (then, \texttt{\textbackslash{}printProofs} only input that
%file). You can input that file from other files if you want to separate
%the body from the proof, but it won't work if you have restate theorems.
%You should use the option \texttt{external\ appendix} (starting from
%version 2022/02/04) to restate appropriately the theorem:
%
%\begin{Shaded}
%\begin{Highlighting}[]
%\BuiltInTok{\textbackslash{}usepackage}\NormalTok{[createShortEnv,conf=\{external appendix\}]\{}\ExtensionTok{proof{-}at{-}the{-}end}\NormalTok{\}}
%\end{Highlighting}
%\end{Shaded}
%
%Then, compile your main file and create a new file for the appendix,
%load the theorems environments as for the main file, in the preambule
%load the \texttt{xr} package to properly have references (otherwise you
%won't have the appropriate number for the theorems) using something
%like:
%
%\begin{verbatim}
%\usepackage{xr}
%\externaldocument{name_of_main_file_without_extension}
%\end{verbatim}
%
%and in this new document just add the appendix wherever you want:
%
%\begin{Shaded}
%\begin{Highlighting}[]
%\FunctionTok{\textbackslash{}includeExternalAppendix}\NormalTok{\{name\_of\_main\_file\_without\_extension\}}
%\end{Highlighting}
%\end{Shaded}
%
%You can see an exemple in the github repository, with the two files
%\texttt{demo\_external\_appendix.tex} and
%\texttt{demo\_external\_appendix\_part2.tex}. Whenever your proofs
%change in the main paper, make sure to compile first the main document
%(so that the auxiliary file gets updated), and then compile the appendix
%document. Note that you can still use
%\texttt{\textbackslash{}printProofs} in the main document if you like
%(while writting your paper, it is certainly easier to keep the proofs
%and theorems in the same document).
%
%\section{List of options}\label{list-of-options}
%
%Here is the list of fundamental options supported. Most options have a
%\texttt{no} version, with \texttt{no} written before. Note that you may
%prefer to use directly the alias/styles (see next paragraph) as the
%options listed here are quite fundamental and atomic.
%
%\begin{itemize}
%\item
%  \texttt{category}: category of the proof (if you want to put proofs at
%  several places), can be anything
%\item
%  \texttt{proof\ here}/\texttt{no\ proof\ here}: put (or not) the proof
%  in the main text
%\item
%  \texttt{proof\ end}/\texttt{no\ proof\ end}: display the proof in
%  appendix
%\item
%  \texttt{restate}/\texttt{no\ restate}: restate the theorem in appendix
%\item
%  \texttt{link\ to\ proof}/\texttt{no\ link\ to\ proof}: Display a link
%  to the proof in the main text
%\item
%  \texttt{opt\ all\ end}/\texttt{no\ opt\ all\ end}: put the theorem and
%  proof only in appendix. You may prefer the alias \texttt{all\ end},
%  that also makes sure that the proof is indeed displayed in appendix.
%\item
%  \texttt{text\ link}: text of the link to the proof, defaults to
%
%  \texttt{\{See\ \textbackslash{}hyperref{[}proof:prAtEnd\textbackslash{}pratendcountercurrent{]}\{proof\}\ on\ page\textasciitilde{}\textbackslash{}pageref\{proof:prAtEnd\textbackslash{}pratendcountercurrent\}.\}}
%\item
%  \texttt{text\ proof}: text displayed in place of ``Proof'' in the
%  appendix. Defaults to
%  \texttt{\{Proof\ of\ \textbackslash{}string\textbackslash{}pratendRef\{thm:prAtEnd\textbackslash{}pratendcountercurrent\}\}}
%\item
%  \texttt{restate\ command}: name of a unique macro (without backslash)
%  that will be defined as an alias to restate the theorem wherever you
%  want
%\item
%  \texttt{restated\ before}: if the theorems has been stated before
%  (with \texttt{\textbackslash{}theoremProofEndRestateBefore}), then we
%  just need to put the restate command in place of the theorem, and
%  enable this option
%\item
%  \texttt{both}/\texttt{no\ both}: only for
%  \texttt{\textbackslash{}textInAppendix}, specifies that the text must
%  be present in both the main text and the appendix.
%\item
%  \texttt{external\ appendix}: to ensure the appendix can be included in
%  another file (see details above), starting from version 2022/02/04.
%\end{itemize}
%
%Here are all the alias/styles (you can create you own as well), they are
%practical to quickly define a behaviours, but are made of the basic
%options listed above:
%
%\begin{itemize}
%\tightlist
%\item
%  \texttt{normal}: like a `normal' theorem, without any proof in the
%  appendix, and with a proof displayed in the main text. Shortcut for
%  \texttt{proof\ here,\ no\ all\ end,\ no\ proof\ end,\ no\ link\ to\ proof,\ no\ restate,\ no\ both}.
%\item
%  \texttt{end}: theorems whose proof need to go in the appendix. Shorcut
%  for \texttt{proof\ at\ the\ end,\ link\ to\ proof}.
%\item
%  \texttt{all\ end}: makes sure both the theorem and the proof are in
%  appendix. Alias of \texttt{end,\ opt\ all\ end}.
%\item
%  \texttt{proof\ at\ the\ end}: theorems whose proof need to go in the
%  appendix contrary to \texttt{end} it does not make sure that there is
%  a link to the proof. Shorcut for
%  \texttt{no\ proof\ here,\ no\ all\ end,\ proof\ end,\ no\ both}.
%\item
%  \texttt{debug}: make sure the proof is written in the main text as
%  well (alias of \texttt{proof\ here,\ no\ opt\ all\ end}), it is quite
%  practical to use when you write a proof to be able to use synctex
%  features to move between the pdf and the file.
%\item
%  \texttt{no\ link\ to\ theorem}: Remove the link from the proof to the
%  theorem, alias of \texttt{text\ proof=\{\textbackslash{}proofname\}}
%\item
%  \texttt{stared} (or \texttt{no\ number}): when you use the stared
%  version of a theorem you don't have any number, so cref fails to write
%  a nice link to the theorem. This option changes the text of ``Proof'',
%  by keeping the link but writting only \texttt{Proof}. Equivalent to
%  \texttt{text\ proof=\{\textbackslash{}string\textbackslash{}mbox\{\textbackslash{}string\textbackslash{}hyperref{[}thm:prAtEnd\textbackslash{}pratendcountercurrent{]}\{\textbackslash{}proofname\}\}\}}
%\item
%  \texttt{see\ full\ proof}: useful when you want to write in the main
%  text only a sketch of proof, this alias writes a link
%  \texttt{See\ full\ proof\ on\ page\ X.}. Equivalent to
%  \texttt{text\ link=\{See\ \textbackslash{}hyperref{[}proof:prAtEnd\textbackslash{}pratendcountercurrent{]}\{full\ proof\}\ on\ page\textasciitilde{}\textbackslash{}pageref\{proof:prAtEnd\textbackslash{}pratendcountercurrent\}.\}}
%\item
%  \texttt{one\ big\ link}: instead of two links, one for page, one for
%  proof, put just one link around everything. It can also accept an
%  optional argument which will be the text of the link, like
%  \texttt{one\ big\ link=Go\ to\ the\ proof}. The default value is
%  \texttt{See\ proof\ on\ page\textasciitilde{}\textbackslash{}pageref*\{proof:prAtEnd\textbackslash{}pratendcountercurrent.\}}.
%\item
%  \texttt{one\ big\ link\ translated}: This is like
%  \texttt{one\ big\ link}, but automatically add the page at the end
%  (and a big link around). Practical to quickly define a translation
%  like \texttt{one\ big\ link\ translated=Voir\ preuve\ page}. See also
%  \texttt{text\ proof\ translated}.
%\item
%  \texttt{text\ link\ section}: Put a link to the proof looking like
%  ``See proof in section X''. Defaults to
%  \texttt{text\ link=\{See\ \textbackslash{}hyperref{[}proof:prAtEnd\textbackslash{}pratendcountercurrent{]}\{proof\}\ in\ \textbackslash{}pratendSectionlikeCref.\}}
%\item
%  \texttt{text\ link\ section\ full\ proof}: Put a link to the proof
%  looking like ``See full proof in section X''. Defaults to
%  \texttt{text\ link=\{See\ \textbackslash{}hyperref{[}proof:prAtEnd\textbackslash{}pratendcountercurrent{]}\{full\ proof\}\ in\ \textbackslash{}pratendSectionlikeCref.\}}
%\item
%  \texttt{default\ text\ link}: default text for the link to the proof,
%  equivalent of
%  \texttt{text\ link=\{See\ \textbackslash{}hyperref{[}proof:prAtEnd\textbackslash{}pratendcountercurrent{]}\{proof\}\ on\ page\textasciitilde{}\textbackslash{}pageref\{proof:prAtEnd\textbackslash{}pratendcountercurrent\}.\}}
%\item
%  \texttt{default\ text\ proof}: default text for the proof in appendix,
%  equivalent of
%  \texttt{text\ proof=\{Proof\ of\ \textbackslash{}string\textbackslash{}pratendRef\{thm:prAtEnd\textbackslash{}pratendcountercurrent\}\}}
%\item
%  \texttt{text\ proof\ only\ theorem}: equivalent of
%  \texttt{text\ proof=\{\textbackslash{}string\textbackslash{}pratendRef\{thm:prAtEnd\textbackslash{}pratendcountercurrent\}\}},
%  helpful if the class uses a custom proof environment that would write
%  \texttt{Proof\ (Proof\ of\ Lemma\ X)} instead of
%  \texttt{Proof\ of\ Lemma\ X.}: adding this style turns it into
%  \texttt{Proof\ (Lemma\ X)}. See more details in the troubleshooting
%  section.
%\item
%  \texttt{text\ proof\ translated}: like \texttt{default\ text\ proof},
%  but takes one argument and use it instead of \texttt{Proof\ of}.
%  Example: \texttt{text\ proof\ translated=\{Preuve\ du\}}
%\item
%  \texttt{text\ link\ external\ appendix}: style used to configure the
%  text to display when the appendix is in an external file. Defaults to
%  \texttt{text\ link=\{The\ proof\ is\ in\ the\ appendix.\}}.
%\item
%  \texttt{external\ appendix}: useful when the the proofs are moved in
%  an external file (read more above). In particular, it will ensure that
%  normal theorems are used instead of restate since restate does not
%  work across files. This will also apply
%  \texttt{text\ link\ external\ appendix} to change the text
%  accordingly.
%\item
%  \texttt{only\ external\ appendix}: like \texttt{external\ appendix},
%  but without changing the text.
%\item
%  \texttt{bare\ defaults}: default style that is loaded before anything
%  else that configure by default a link to the proof, put the proof in
%  appendix, use the category \texttt{defaultcategory}. It is an alias of
%  \texttt{end,\ link\ to\ proof,\ no\ restate,category=defaultcategory,\ default\ text\ link,default\ text\ proof,restate\ command=pratenddummymacro}.
%\item
%  \texttt{configuration\ options}: style that contains the options used
%  to load the package. It is called right after \texttt{bare\ defaults}.
%  Note that you cannot insert macro in the options, overwrite
%  \texttt{global\ custom\ defaults} instead
%\item
%  \texttt{global\ custom\ defaults}: empty style that you can overwrite
%  to change the global defaults
%\item
%  \texttt{local\ custom\ defaults}: empty style that you can overwrite
%  to change the ``local'' defaults, like category
%\item
%  \texttt{all\ defaults}: all the defaults, equivalent of
%  \texttt{bare\ defaults,\ configuration\ options,\ global\ custom\ defaults,\ local\ custom\ defaults}
%\end{itemize}
%
%\section{Package options}\label{package-options}
%
%The package comes with multiple options: - The
%\texttt{conf=\{CONFIGURATION\}} option is used to configure the default
%style. Usage:
%\texttt{\textbackslash{}usepackage{[}conf=\{normal,\ text\ link\ section\}{]}\{proof-at-the-end\}}
%- The \texttt{disablePatchSection} is useful to stop the library from
%patching the section-like commands and chapters (by default, we
%automatically add a \texttt{\textbackslash{}label} at the end of the
%section which is used by \texttt{text\ link\ section} to find the
%Appendix). If you enable this option, you can manually add the label by
%using \texttt{\textbackslash{}pratendAddLabel} after the appendix
%section (without any argument). - \texttt{commandRef=NAMECOMMAND}: By
%default, theorem are referenced to using
%\texttt{\textbackslash{}autoref\{\}}. You can change the value of
%\texttt{NAMECOMMAND} to use another command, like
%\texttt{commandRef=Cref} to use \texttt{\textbackslash{}Cref\{\}}. In
%that case, make sure to load both \texttt{hyperref} and
%\texttt{cleveref} (typically at the very end of your preambule) as
%\texttt{cleveref} needs to be loaded after \texttt{hyperref}.
%Internally, the library creates and uses
%\texttt{\textbackslash{}pratendRef} instead of
%\texttt{\textbackslash{}autoref}. - \texttt{createShortEnv}: creates
%(using \texttt{newEndThm} and \texttt{newEndProof}) the shortcut
%environments \texttt{proofE} (to replace \texttt{proof}), \texttt{thmE}
%(to replace \texttt{thm}), \texttt{theoremE} (to replace\ldots{} ok you
%got the pattern), \texttt{lemmaE} and \texttt{corollaryE}. Note that it
%is still your job to define the \texttt{thm}, \texttt{lemma}
%environments.
%
%\section{Troubleshooting}\label{troubleshooting}
%
%Here are some common issues you may have, with explainations to solve
%them.
%
%\subsection{The link after the proof does not mention the section, but
%something else (an equation, a figure,
%etc.)}\label{the-link-after-the-proof-does-not-mention-the-section-but-something-else-an-equation-a-figure-etc.}
%
%Sometimes, you may see something like
%\texttt{See\ proof\ in\ Equation\ A} instead of
%`\texttt{See\ proof\ in\ Appendix\ A}. This issue
%\href{https://github.com/leo-colisson/proof-at-the-end/issues/2}{was
%reported before} and should be solved on newer versions (starting from
%2022/01/27) by automatically patching sections. If you have this issue,
%just upgrade (for instance by copying the
%\href{https://github.com/leo-colisson/proof-at-the-end}{\texttt{proof-at-the-end.sty}
%file} at the root of your project). We proceed by patching the
%\texttt{\textbackslash{}chapter}, \texttt{\textbackslash{}section},
%\texttt{\textbackslash{}subsection}, \texttt{subsubsection} and
%\texttt{\textbackslash{}paragraph} commands to add
%\texttt{\textbackslash{}pratendAddLabel} which will help the package to
%find the label of the current section. If you prefer the old behavior,
%you can disable it using the \texttt{disablePatchSection} option (more
%details in the package options).
%
%\subsection{\texorpdfstring{I get an error
%\texttt{ERROR:\ File\ ended\ while\ scanning\ use\ of\ \textbackslash{}@xverbatim.}
%when using verbatim inside a
%theorem}{I get an error ERROR: File ended while scanning use of \textbackslash@xverbatim. when using verbatim inside a theorem}}\label{i-get-an-error-error-file-ended-while-scanning-use-of-xverbatim.-when-using-verbatim-inside-a-theorem}
%
%I turns out that this library needs to use the \texttt{+b} option of
%environments in order to manipulate appropriately the theorem/proof.
%Unfortunately, it means that it is not possible to use verbatim
%environments inside. While
%\href{https://tex.stackexchange.com/questions/489435/use-environment-into-new-xparse-environment}{this
%answer} suggests that it is impossible to avoid this issue, other people
%reported that it may be possible to use catcodes to capture the
%environment body verbatim, before using scantokens to reparse it, or to
%use
%\href{https://mirror.ibcp.fr/pub/CTAN/macros/latex/contrib/filecontentsdef/filecontentsdef.pdf}{filecontentsdef}
%(see for instance
%\href{https://tex.stackexchange.com/questions/631810/create-a-newdocumentenvironment-programmatically?noredirect=1\#comment1575605_631810}{this
%comment}{]}, ). However, it may be possible that this make syntex
%unusable, so it may not worth the effort.
%
%In anycase, there exists some workarounds, some of the are for instance
%give in the \href{https://texfaq.org/FAQ-verbwithin}{TeX FAQ}. The
%simplest solution, if your code is simple/short enough, is to use
%\texttt{\textbackslash{}texttt} (you can replace backslash with
%\texttt{\textbackslash{}textbackslash} inside etc, and include it in a
%\texttt{verse} if you have multiple lines). If you have a longer text,
%you may also like \texttt{lrbox} to put your content inside a box, for
%instance like that:
%
%\begin{Shaded}
%\begin{Highlighting}[]
%\CommentTok{\%\% Create the box}
%\FunctionTok{\textbackslash{}newsavebox}\NormalTok{\{}\FunctionTok{\textbackslash{}myEndBox}\NormalTok{\}}
%
%\KeywordTok{\textbackslash{}begin}\NormalTok{\{}\ExtensionTok{lrbox}\NormalTok{\}\{}\FunctionTok{\textbackslash{}myEndBox}\NormalTok{\}}
%  \KeywordTok{\textbackslash{}begin}\NormalTok{\{}\ExtensionTok{minipage}\NormalTok{\}\{1.0}\FunctionTok{\textbackslash{}linewidth}\NormalTok{\}}
%\KeywordTok{\textbackslash{}begin}\NormalTok{\{}\ExtensionTok{verbatim}\NormalTok{\}}
%\VerbatimStringTok{\textbackslash{}newEndThm[normal]\{thmE\}\{thm\}}
%\VerbatimStringTok{\textbackslash{}newEndProof\{proofE\}\{proof\}}
%\KeywordTok{\textbackslash{}end}\NormalTok{\{}\ExtensionTok{verbatim}\NormalTok{\}}
%  \KeywordTok{\textbackslash{}end}\NormalTok{\{}\ExtensionTok{minipage}\NormalTok{\}}
%\KeywordTok{\textbackslash{}end}\NormalTok{\{}\ExtensionTok{lrbox}\NormalTok{\}}
%
%\CommentTok{\%\% Create a theorem:}
%\KeywordTok{\textbackslash{}begin}\NormalTok{\{}\ExtensionTok{thmE}\NormalTok{\}}
%\NormalTok{  Insert the box:}
%  
%  \KeywordTok{\textbackslash{}begin}\NormalTok{\{}\ExtensionTok{center}\NormalTok{\}}
%    \FunctionTok{\textbackslash{}usebox}\NormalTok{\{}\FunctionTok{\textbackslash{}myEndBox}\NormalTok{\}}
%  \KeywordTok{\textbackslash{}end}\NormalTok{\{}\ExtensionTok{center}\NormalTok{\}}
%\KeywordTok{\textbackslash{}end}\NormalTok{\{}\ExtensionTok{thmE}\NormalTok{\}}
%\end{Highlighting}
%\end{Shaded}
%
%\subsection{\texorpdfstring{The proof writes
%\texttt{Proof\ (Proof\ of\ Lemma\ X)} (or something else) instead of
%\texttt{Proof\ of\ Lemma\ X.}}{The proof writes Proof (Proof of Lemma X) (or something else) instead of Proof of Lemma X.}}\label{the-proof-writes-proof-proof-of-lemma-x-or-something-else-instead-of-proof-of-lemma-x.}
%
%If the proof starts with something different from
%\texttt{Proof\ of\ Lemma\ X}, you are likely to use a class/package that
%defines a \texttt{proof} environment different from the one defined in
%\texttt{amsthm}, like \texttt{svjour3} by Springer. To fix this, you
%have two options: if you want to keep using the original \texttt{proof}
%environment, and get something like \texttt{Proof\ (Lemma\ X)}, you can
%either configure:
%
%\begin{verbatim}
%\pratendSetGlobal{
%  text proof={\string\pratendRef{thm:prAtEnd\pratendcountercurrent}},
%}
%\end{verbatim}
%
%or, if you use a version greater than 2025/06/11 (included), you can
%simply call the \texttt{text\ proof\ only\ theorem}, either like in
%\texttt{\textbackslash{}pratendSetGlobal\{text\ proof\ only\ theorem\}}
%or directly when loading the package:
%
%\begin{verbatim}
%\usepackage[createShortEnv, conf={text proof only theorem}]{proof-at-the-end}
%\end{verbatim}
%
%Another option is to replace your class \texttt{proof} environment with,
%for instance, the one provided by \texttt{amsthm} like in:
%
%\begin{verbatim}
%% Removes the environment 'proof' defined by the class and replaces it with the one of amsthm:
%\let\proof\undefined
%\let\endproof\undefined
%\usepackage{amsthm}
%\end{verbatim}
%
%but ensure that \texttt{amsthm} does not produce additional conflicts.
%This method will write \texttt{Proof\ of\ Lemma\ X.} before starting the
%proof.
%
%\section{Changelog}\label{changelog}
%
%\begin{itemize}
%\tightlist
%\item
%  2026/01/12: Fix a bug that was not allowing stuff like
%  \texttt{\textbackslash{}cite{[}foo{]}\{bar\}} in theorem titles by
%  introducing an extra group around the title
%\item
%  2025/06/11: Add \texttt{text\ proof\ only\ theorem}
%\item
%  2022/02/07:
%
%  \begin{enumerate}
%  \def\labelenumi{\arabic{enumi}.}
%  \tightlist
%  \item
%    \href{https://github.com/leo-colisson/proof-at-the-end/issues/2}{issue
%    2} was not really solved in fact\ldots{} Now, I use a more robust
%    method to detect the current section: I write in an AUX file the
%    label of the proof section. This way, I don't need anymore to use
%    dirty tricks to recover the section label.
%  \item
%    Change the default text when the appendix is in an external file,
%    and add options like \texttt{text\ link\ external\ appendix} to
%    customize the text appearing when the appendix is in an external
%    file. Of course, you can still change it using \texttt{text\ link}
%    directly.
%  \end{enumerate}
%\item
%  2022/02/04:
%
%  \begin{enumerate}
%  \def\labelenumi{\arabic{enumi}.}
%  \tightlist
%  \item
%    Add a way to put theorems in different files. (see
%    \texttt{external\ appendix})
%  \item
%    Change the path for auxiliary files (should be transparent for the
%    user)
%  \item
%    Add commands to change local/global configuration easily (see
%    \texttt{\textbackslash{}pratendSetLocal},\texttt{\textbackslash{}pratendSetGlobal})
%  \end{enumerate}
%\item
%  2022/02/01:
%
%  \begin{enumerate}
%  \def\labelenumi{\arabic{enumi}.}
%  \tightlist
%  \item
%    Fix a typo when defining the shortcut for lemma
%  \item
%    Add a shortcut for proposition
%  \item
%    Fix a small spacing issue in the text before the proof in Appendix
%  \end{enumerate}
%\item
%  2022/01/28:
%
%  \begin{enumerate}
%  \def\labelenumi{\arabic{enumi}.}
%  \tightlist
%  \item
%    Fix the issue when using sharps in a proof
%    https://github.com/leo-colisson/proof-at-the-end/issues/7.
%  \item
%    Provide \texttt{\textbackslash{}newEndThm},
%    \texttt{\textbackslash{}newEndProof} and the option
%    \texttt{createShortEnv} to quickly create environments.
%  \end{enumerate}
%\item
%  2022/01/27:
%
%  \begin{enumerate}
%  \def\labelenumi{\arabic{enumi}.}
%  \tightlist
%  \item
%    Patch chapters/sections/\ldots{} to better detect the current
%    section and fix
%    \href{https://github.com/leo-colisson/proof-at-the-end/issues/2}{issue
%    2}. This can be disabled using the package option
%    \texttt{disablePatchSection}.
%  \item
%    Add an option \texttt{commandRef} to use other ref libraries, like
%    \texttt{cleveref} instead of \texttt{autoref}.
%  \item
%    Normal restatable theorem do not need to be restated. Fix
%    \href{https://github.com/leo-colisson/proof-at-the-end/issues/8}{issue
%    8}.
%  \end{enumerate}
%\end{itemize}
%
%\section{Contributions}\label{contributions}
%
%Feel free to contribute, report bugs, and send pull requests on the
%github repository \url{https://github.com/leo-colisson/proof-at-the-end}
%!
%
%NB: the documentation is generated from the Markdown file
%\texttt{README.md} thanks to pandoc. These commands may help you:
%
%\begin{Shaded}
%\begin{Highlighting}[]
%\ExtensionTok{\%\%}\NormalTok{ Compile the demo}
%\FunctionTok{make}\NormalTok{ demo}
%\ExtensionTok{\%\%}\NormalTok{ Clean the project}
%\FunctionTok{make}\NormalTok{ clean}
%\ExtensionTok{\%\%}\NormalTok{ Generate the documentation}
%\FunctionTok{make}\NormalTok{ doc}
%\ExtensionTok{\%\%}\NormalTok{ Generate a package for CTAN}
%\FunctionTok{make}\NormalTok{ package}
%\end{Highlighting}
%\end{Shaded}
%
%
%
%
%\StopEventually{}
%\section{The Code}
%\iffalse
%    \begin{macrocode}
%<*proof-at-the-end.sty>
%    \end{macrocode}
%\fi
\NeedsTeXFormat{LaTeX2e}
\ProvidesPackage{proof-at-the-end}[2026/01/12 A package to move proofs in appendix]
 
\RequirePackage{etoolbox}
\RequirePackage{thmtools}
\RequirePackage{thm-restate}
\RequirePackage{catchfile}
\RequirePackage{pgfkeys}
\RequirePackage{xparse}
% Load hyperref at the end
\AtEndPreamble{%
  \RequirePackage{hyperref}
}%
\RequirePackage{kvoptions}

%% https://tex.stackexchange.com/questions/109747/put-all-package-options-into-one-command
%% Forward the options list to the command \pratendOptdefconf
%% in order to use:
%% \usepackage[conf={normal}]{proof-at-the-end}
\SetupKeyvalOptions{
  family=pratendOpt,
  prefix=pratendOpt, % "Namespace": Prefix of the commands, conditions....
}
\DeclareStringOption{conf}
\DeclareLocalOptions{conf}% \pratendOptdefconf contains the proof
\DeclareBoolOption{disablePatchSection} % Use \ifpratendOptdisablePatchSection ... \else ... \fi
\DeclareStringOption[autoref]{commandRef} % Use commandRef=cref to use cleveref. Defaults to autoref.
\DeclareBoolOption{createShortEnv} % To create thmE, theoremE, proofE, lemmaE, corrolaryE
\ProcessKeyvalOptions*

% \DeclareOption*{\PackageWarning{proof-at-the-end}{Unknown}}
% \ProcessOptions\relax

%% https://tex.stackexchange.com/a/13289/116348
\newwrite\appendwrite
\newcount\pratend@count@makeallother

%%% Loop to make sure all chars are normal letters, including #.
\newcommand{\makeallother}{%
  \pratend@count@makeallother0\relax
  \loop\ifnum\pratend@count@makeallother<255\relax
  \catcode\pratend@count@makeallother12\relax
  \advance\pratend@count@makeallother by 1\relax
  \repeat
}

%%% This command allows the user to append things to a file. But the content may contain macros that
%%% will be evaluated before writing them to the file.
%%% TODO: Efficiency may be improved by writing to the file only before the printProof.
% The first argument is the file name
% The second argument is the text to write
\NewDocumentCommand\appendtofile{m+m}{%
  \begingroup
  %% Read the content
  \begingroup%
  \IfFileExists{#1}{%
    \newlinechar\endlinechar%
    \makeallother% Turn all chars into normal letters.
    \everyeof{\noexpand}%
    \edef\fileContent{\@@input #1 }%
  }{%
    \let\fileContent\empty
  }%
  %% Open the file to write in it:
  \immediate\openout\appendwrite #1\relax%
  %% Write the text
  \immediate\write\appendwrite{\fileContent}%
  \endgroup%
  \immediate\write\appendwrite{#2}%
  %% Close the file
  \immediate\closeout\appendwrite%
  \endgroup
}

%% Get the prefix of the file, given the jobname.
\NewDocumentCommand\pratendGeneratePrefixFile{m}{#1-pratend}

\def\prefixPrAtEndFiles{\pratendGeneratePrefixFile{\jobname}}


%% This functions takes one input: the category (without .tex),
%% If it's the first time we write in this
%% category file, it "blanks" it.

\newcommand*\eraseIfNeeded[1]{%
  % A macro 'pratendmacrocat{category}' is created to check
  % if it's the first time we write in this category file.
  \protected@edef\macroname{pratendmacrocat#1}%
  \ifcsdef{\macroname}{% The macro exists, nothing to do
  }{ % The macro does not exists, create it, and empty the file
    \global\expandafter\def\csname \macroname\endcsname{true}%
    \immediate\openout\appendwrite=\prefixPrAtEndFiles#1.tex%
    \immediate\write\appendwrite{}%
    \immediate\closeout\appendwrite%
  }%
}

\newif\ifproofhere
\newif\ifproofend
\newif\ifrestatethm
\newif\ifexternalAppendix% True if the output file should me inputable in another file (using xr for crossref). Useful to remote restate.
\newif\iflinktoproof
\newif\ifboth
\newif\ifallattheend
\newif\ifrestatedbefore
\pgfkeys{
  /prAtEnd/.cd, %% Proof at end will be the main path
  %% Category of the proof (if you want to put proofs
  %% at several places), can be anything
  category/.initial=defaultcategory,
  category/.store in=\category,
  category/.get=\category,
  %% Display the proof in the main part
  proof here/.is if=proofhere,
  no proof here/.style={proof here=false}, % alias
  %% Display the proof when using \printProofs
  proof end/.is if=proofend,
  no proof end/.style={proof end=false},
  %% Restate the theorem when using \printProofs
  only external appendix/.is if=externalAppendix,
  text link external appendix/.style={
    text link={The proof is in the appendix.}
  },
  external appendix/.style={
    only external appendix,
    text link external appendix,
  },
  %% Restate the theorem when using \printProofs
  restate/.is if=restatethm,
  no restate/.style={restate=false},
  %% Put a link to the proof after the theorem
  link to proof/.is if=linktoproof,
  no link to proof/.style={link to proof=false},
  %% Put the theorem and proof only in appendix
  opt all end/.is if=allattheend,
  no opt all end/.style={opt all end=false},
  %% Text of link
  text link/.code={\def\pratendtextlink{#1}},
  %% Text of proof. Make sure also to "\renewcommand*{\proofname}{Name of the proof}"
  %% to make sure the proof for normal theorems are changed
  text proof/.code={\def\pratendtextproof{#1}},
  %% Custom restate command
  restate command/.code={\protected@edef\pratendcustomrestate{#1}},
  %% (Re)stated before
  %% If the theorems has been stated before, then we just need to put the restate command in
  %% place of the argument, and we set this value to true:
  restated before/.is if=restatedbefore,
  no restated before/.style={restated before=false},
  %% In star version, we don't want 
  %% Put the text (defined only for \textInAppendix) in both the
  %% current location and in appendix
  both/.is if=both,
  no both/.style={both=false},
  %%%% Alias and styles
  normal/.style={
    proof here,
    no opt all end,
    no proof end,
    no link to proof,
    no restate,
    no both,
  },
  proof at the end/.style={
    no proof here,
    no opt all end,
    proof end,
    no both,
  },
  end/.style={
    proof at the end,
    link to proof,
  },
  all end/.style={
    end,
    opt all end,
  },
  debug/.style={
    no opt all end,
    proof here
  },
  no link to theorem/.style={ % Remove the link to the theorem
    text proof={\proofname},
  },
  stared/.style={ % Remove 
    text proof={\string\mbox{\string\hyperref[thm:prAtEnd\pratendcountercurrent]{\proofname}}},
  },
  no number/.style={
    stared
  },
  see full proof/.style={
    text link={See \hyperref[proof:prAtEnd\pratendcountercurrent]{full proof} on page~\pageref{proof:prAtEnd\pratendcountercurrent}.}
  },
  one big link/.style={
    text link={\hyperref[proof:prAtEnd\pratendcountercurrent] {#1}}
  },
  one big link/.default={%
    See proof on page~\pageref*{proof:prAtEnd\pratendcountercurrent}.
  },
  one big link translated/.style={
    one big link={#1~\pageref*{proof:prAtEnd\pratendcountercurrent}.}
  },
  text link section/.style={%
      text link={See \hyperref[proof:prAtEnd\pratendcountercurrent]{proof} in \pratendSectionlikeCref.}
  },
  text link section full proof/.style={%
      text link={See \hyperref[proof:prAtEnd\pratendcountercurrent]{full proof} in \pratendSectionlikeCref.}
  },
  default text link/.style={
    text link={See \hyperref[proof:prAtEnd\pratendcountercurrent]{proof} on page~\pageref{proof:prAtEnd\pratendcountercurrent}.},
  },
  text proof translated/.style={
    text proof={#1 \string\pratendRef{thm:prAtEnd\pratendcountercurrent}},
  },
  default text proof/.style={
    text proof={Proof of \string\pratendRef{thm:prAtEnd\pratendcountercurrent}},
  },
  text proof only theorem/.style={
    text proof={\string\pratendRef{thm:prAtEnd\pratendcountercurrent}},
  },
  %%%% Defaults
  bare defaults/.style={
    end,
    link to proof,
    no restate,
    category=defaultcategory,
    default text link,
    default text proof,
    restate command=pratenddummymacro,
  },
  configuration options/.style/.expand once={
    % This styles will contain the configuration
    % given as options of the package like:
    % \usepackage[conf={normal, no link to proof}]{proof-at-the-end}
    % The package options does not accept macros and valued keys
    % due to some fundamental issues:
    % https://tex.stackexchange.com/questions/489564/use-unexpanded-macro-in-package-options/489570#489570
    % so if you need to write macro/valued key, edit instead the
    % style "/prAtEnd/global custom defaults" or
    % "local /prAtEnd/custom defaults" for local changes instead.
    \pratendOptconf%
  },
  global custom defaults/.style={
    %% you can put in this style any global defaults
    %% that should overwrite the usual defaults.
  },
  local custom defaults/.style={
    %% you can put in this style any overwrite of the defaults
    %% that should be "local" and changed over the file, like
    % the category for a given section.
  },
  all defaults/.style={
    %% Load all the style that sets the default values
    bare defaults,
    configuration options,
    global custom defaults,
    local custom defaults,
  },
}

\newcounter{counterAllProofEnd}
\stepcounter{counterAllProofEnd}

\NewDocumentEnvironment{theoremEndRestateBefore}{mO{}m+b}{
  \stepcounter{counterAllProofEnd}%
  \protected@edef\currcounterval{\roman{counterAllProofEnd}}%
  \expandafter\protected@xdef\csname #3\endcsname{\currcounterval}%
  \expanded{%
    \noexpand\begin{restatable*}[{\unexpanded{#2}}]{\unexpanded{#1}}{prAtEndRestate\currcounterval}\noexpand\label{thm:prAtEnd\currcounterval}%
    \unexpanded{#4}%
    \noexpand\end{restatable*}%
  }%
}{}

\NewDocumentEnvironment{theoremEndOrig}{O{}mO{}+b}{
  % The first facultative argument will be the options: type of proof you want, the file to which you want to write...
  % The first mandatory option is the type of the theorem (thm,lemma,...)
  % The second facultative argument will be the title
  % the second mandatory option is the theorem (will \label inside eventually)
  % the last mandatory option is the proof, or, if restated, the name of the theorem.
  \global\def\pratendlastoptions{#1}%
  \pgfkeys{%
    /prAtEnd/.cd,
    all defaults,
    #1
  }%
  \stepcounter{counterAllProofEnd}%
  %% Create the file if it's the first time
  \eraseIfNeeded{\category}%
  %% If the theorem must be written here:
  \unless\ifallattheend%
    %% Restate the theorem if it was stated before:
    \ifrestatedbefore%
      \protected@xdef\pratendcountercurrent{\csname #4\endcsname} % Store the current (alpha value of the) counter
      \csname prAtEndRestate\pratendcountercurrent\endcsname % Restate the theorem
    \fi%
    %%  Otherwise just state the theorem in a restatable environment
    \unless\ifrestatedbefore%
      \expanded{% This {} ----------V is needed to allow stuff like \cite[foo]{bar}. Test also with $\mathtt{G}$
        \noexpand\begin{restatable}[{\unexpanded{#3}}]{\unexpanded{#2}}{prAtEndRestate\roman{counterAllProofEnd}}%
          \noexpand\label{thm:prAtEnd\roman{counterAllProofEnd}}%
          \unexpanded{#4}%
          \noexpand\end{restatable}%
      }%
      %% Store the current (alpha value of the) counter
      %% in \pratendcountercurrent
      \protected@xdef\pratendcountercurrent{\roman{counterAllProofEnd}} %
    \fi%
    %% Create a custom alias to restate the theorem
    \expandafter\protected@xdef\csname \pratendcustomrestate\endcsname{\noexpand\csname prAtEndRestate\pratendcountercurrent\endcsname}%
    %% Restate the theorem if needed in appendix.
    %% As explained in https://github.com/leo-colisson/proof-at-the-end/issues/8
    %% No need to restate of there is no proof at the end.
    \ifproofend%
      \ifrestatethm%
        \ifexternalAppendix% We write the full theorem, useful to compile in a separate file (use xr for references)
          \appendtofile{\prefixPrAtEndFiles\category.tex}{\string\begingroup\string\renewcommand\string\the\detokenize{#2}{\string\ref{thm:prAtEnd\pratendcountercurrent}}\string\begin{\detokenize{#2}}[#3]\detokenize{#4}\string\end{\detokenize{#2}}\string\endgroup}%
        \else%
          \appendtofile{\prefixPrAtEndFiles\category.tex}{\string\prAtEndRestate\pratendcountercurrent*}%
        \fi%
      \fi%
    \fi%
  \fi%  
  %% If the theorem is not stated in the main text,
  %% write it at the end 
  \ifallattheend%
    %% Store the current (alpha value of the) counter
    %% in \pratendcountercurrent
    \protected@xdef\pratendcountercurrent{\roman{counterAllProofEnd}}
    %% Create a custom alias to restate the theorem
    \expandafter\protected@xdef\csname \pratendcustomrestate\endcsname{\noexpand\csname prAtEndRestate\pratendcountercurrent\endcsname}%
    \appendtofile{\prefixPrAtEndFiles\category.tex}{\string\begin{restatable}[\detokenize{#3}]{#2}{prAtEndRestate\pratendcountercurrent}\string\label{thm:prAtEnd\pratendcountercurrent}\detokenize{#4}\string\end{restatable}}%
  \fi%
}{}


%% See proofEnableDebugSynctex
%%% WARNING: if you change this function, change also proofDisableDebugSynctex
\NewDocumentEnvironment{theoremEnd}{O{}mO{}+b}{%
  \begin{theoremEndOrig}[#1]{#2}[#3]
    #4%
  \end{theoremEndOrig}%
}{}

\NewDocumentEnvironment{proofEndOrig}{O{}+b}{%
  \pgfkeys{%
    /prAtEnd/.cd,
    all defaults,
    prAtEndTmpStyle/.style/.expand once={\pratendlastoptions},
    prAtEndTmpStyle,
    #1
  }%
  \unless\ifallattheend%
    %% Write eventually a link to the proof
    \iflinktoproof%
      \pratendtextlink{}%
    \fi%
    %% And eventually the proof
    \ifproofhere%
      \begin{proof}%
        #2%
      \end{proof}%
    \fi%
  \fi%
  %% Write the proof at the end
  \ifproofend%
  \appendtofile{\prefixPrAtEndFiles\category.tex}{%
    %% The label alone adds some space, see e.g.
    %% https://tex.stackexchange.com/questions/600118/selectlanguage-in-the-middle-of-text-causes-extra-vertical-space/600142#600142
    %% and https://tex.stackexchange.com/questions/600577/adding-a-label-between-two-proofs-adds-unwanted-space
    %% The solution is to use \\Hy@SaveLastskip and \Hy@RestoreLastskip from hyperref
    %% See also
    %% https://tex.stackexchange.com/questions/553394/reliably-add-a-label-to-the-current-section-even-if-the-label-is-not-placed-rig/600847#600847
    \string\makeatletter\string\Hy@SaveLastskip%
    %% This tries to recover the label of the current section... But it is very bad. I keep it only for backward compatibility. I tried to fix it by changing the value of \@counter, but it's not always working
    %% as cref is not using it apparently, and it turns Appendix A into appendix a after using phantomsection.
    %% https://github.com/leo-colisson/proof-at-the-end/issues/2
    \string\label{proofsection:prAtEnd\pratendcountercurrent}% WARNING <-- do not use this label in new documents!
    % the problem is that setting the \@currentlabel does not work for cref, and turns Appendix A
    % into section A when a phantomsection lies in between (not sure why):
    % https://github.com/leo-colisson/proof-at-the-end/issues/2
    %% Now, we prefer a different method: instead we generate a single label per section, and we generate
    %% a new macro in the aux file that will save this label, by creating one such macro per theorem.
    %% This is basically doing something like \pratend@section@for@proofiii{prAtEnd:sectionlike3}
    %% Note that I need sometimes 3 \string because I write two times: the first time I write in the FILE-pratendCATEGORY.tex
    %% file, and this file will be in charge of writting into the .aux file.
    \string\ifdefined\string\pratend@current@sectionlike@label% But before, we check if the section was defined, or the code may fail if no section are defined.
      \string\immediate\string\write\string\@auxout{%
        \string\string\string\gdef\string\string\string\pratend@section@for@proof\pratendcountercurrent{\string\pratend@current@sectionlike@label}%
      }%
    \string\fi
    % See above
    \string\Hy@RestoreLastskip\string\makeatother%
    % We start the proof
    \string\begin{proof}[\pratendtextproof]%
      % We add a label to have a reference to the section
      \string\phantomsection\string\label{proof:prAtEnd\pratendcountercurrent}%
      % We add the proof
      \detokenize{#2}%
      % We end the proof
      \string\end{proof}}%
  \fi%
}{}

%%% WARNING: if you change this function, change also proofDisableDebugSynctex
\NewDocumentEnvironment{proofEnd}{O{}+b}{%
  \begin{proofEndOrig}[#1]%
    #2
  \end{proofEndOrig}%
}{}


%%%%% For synctex debugging
%% Turns https://github.com/leo-colisson/proof-at-the-end/issues/10
\NewDocumentCommand{\pratendEnableDebugSynctex}{}{%
  \RenewDocumentEnvironment{theoremEnd}{O{}mO{}}{\begin{##2}}{\end{##2}}
  \RenewDocumentEnvironment{proofEnd}{O{}}{\begin{proof}}{\end{proof}}
}

%% Turns https://github.com/leo-colisson/proof-at-the-end/issues/10
\NewDocumentCommand{\pratendDisableDebugSynctex}{}{%
  \RenewDocumentEnvironment{theoremEnd}{O{}mO{}+b}{%
    \begin{theoremEndOrig}[##1]{##2}[##3]%
      ##4%
    \end{theoremEndOrig}%
  }{}%
  \RenewDocumentEnvironment{proofEnd}{O{}+b}{%
    \begin{proofEndOrig}[##1]%
      ##2%
    \end{proofEndOrig}%
  }{}
}

\NewDocumentEnvironment{proofEndDebug}{O{}}{%
  \begin{proof}%
  }{%
  \end{proof}%
}
\NewDocumentEnvironment{theoremEndDebug}{O{}mO{}}{\begin{#2}}{\end{#2}}


%%%%% Text in appendix

\NewDocumentEnvironment{textAtEnd}{O{}+b}{
  % Use it to put normal text in Appendix.
  \pgfkeys{
    /prAtEnd/.cd,
    all defaults,
    #1
  }%
  \ifboth%
  #2%
  \fi%
  \eraseIfNeeded{\category}%
  \appendtofile{\prefixPrAtEndFiles\category.tex}{\detokenize{#2}}%
}{}

\NewDocumentCommand\textEnd{O{}+m}{%
  \begin{textAtEnd}[#1]%
    #2%
  \end{textAtEnd}%
}

\NewDocumentCommand\printProofs{O{defaultcategory}}{
  \input{\prefixPrAtEndFiles#1.tex}
}


%%% You can easily modify the defaults:
% \pgfkeys{/prAtEnd/custom defaults/.style={
%     category=greattheorem
%   }
% }
%%% Or create new styles to apply:
% \pgfkeys{/prAtEnd/great category/.style={
%     category=greattheorem
%   }
% }

%%%%%%%%%%%%%%%%%%%%%%%%%%%%%%%%%
%%%   Patch sections, chapters...
%%%%%%%%%%%%%%%%%%%%%%%%%%%%%%%%%
% We patch the \section, \subsection and \chapter commands to make sure that
% the text "See proof in Section XX" points to the appropriate section/subsection (option conf={text link section}).
% We basically maintain a counter counting how many times we saw a section/subsection/... and use it to label
% directly the section afterwards.
\newcounter{pratend@counter@sectionlike}
%% Command to add after a label like "prAtEnd:sectionlike10". Useful to patch \section, \chapter...
\NewDocumentCommand\pratendAddLabel{}{%
  %% First, we get the label
  %% Cool explaination of expandafter, edef, noexpand:
  %%% https://www.tug.org/TUGboat/tb09-1/tb20bechtolsheim.pdf
  %% The label \pratend@current@sectionlike@label is used later.
  \xdef\pratend@current@sectionlike@label{prAtEnd:sectionlike\thepratend@counter@sectionlike}%
  \expandafter\label{\pratend@current@sectionlike@label}%
  \stepcounter{pratend@counter@sectionlike}%
}%
%% Returns the label of the section (in appendix) that contains the proof of the current theorem.
%% The first time, it will return an empty label.
\NewDocumentCommand\pratendLabelProofSection{}{%
  \ifcsname pratend@section@for@proof\pratendcountercurrent\endcsname%
    \csname pratend@section@for@proof\pratendcountercurrent\endcsname%
  \fi%
}%
\NewDocumentCommand\pratendSectionlikeCref{}{%
  % \pratendRef{proofsection:prAtEnd\pratendcountercurrent}% This is not robust enough. Lets change strategy.
  \expandafter\pratendRef{\pratendLabelProofSection}%
}
\ifpratendOptdisablePatchSection%
\else  %
  %% Add the command \pratendAddLabel to \chapter/sections… if the command exists
  %% KOMA class need a special treatment
  \@ifundefined{KOMAClassName}{%
    % For articles/...
    \ifdef{\@chapter}{
      \ifpatchable*{\@chapter}{
        \apptocmd{\@chapter}{\pratendAddLabel}{}{}%
      }{\PackageWarning{proof-at-the-end}{Chapters are not patchable.}{}}
    }{}
    %% patch sections/subsections/paragraph/...
    %% https://tex.stackexchange.com/questions/631713/patch-section-command-fails/631717?noredirect=1#comment1575337_631717
    \ifdef{\@sect}{
      % \apptocmd{\section}{\pratendAddLabel}{}{}%
      \ifpatchable*{\@sect}{
        \apptocmd{\@sect}{\pratendAddLabel}{}{}%
      }{\PackageWarning{proof-at-the-end}{Sections are not patchable.}{}}
    }{}%
  }{%
    %% For KOMA classes/... or it won't work with scrreport
    %% \AddtoDoHook{heading/endgroup/chapter} is too specific 
    \AddtoDoHook{heading/endgroup/chapter}{\pratendAddLabel}%
    \AddtoDoHook{heading/endgroup/section}{\pratendAddLabel}%
  }%
\fi

%%%%%%%%%%%%%%%%%%%%%%%%%%%%%%
%%% Allow cref/autoref/...
%%%%%%%%%%%%%%%%%%%%%%%%%%%%%%
%% Instead of \autoref{proof:blabla} use \pratendRef{proof:blabla}, it allows the ref command to be customized
%% as an option of the package.
\NewDocumentCommand\pratendRef{m}{%
  \csname \pratendOptcommandRef\endcsname{#1}%
}


%%%%%%%%%%%%%%%%%%%%%%%%%%%%%%%%
%%% Create automatically wrapper
%%%%%%%%%%%%%%%%%%%%%%%%%%%%%%%%

%% Defines a shortcut to avoid typing the whole theoremEnd environment.
%% Avoid to manually type \NewDocumentEnvironment...
%% Usage: \newEndThm[default options]{new environment}{old environment}
% E.g.: \newEndThm[normal]{thmE}{thm}{}
\NewDocumentCommand{\newEndThm}{O{}mm}{%
  % \expandafter\NewDocumentEnvironment{#2}{O{}O{}+b}{%
    % \begin{theoremEnd}[#1,##2]{#3}[##1]%
      % ##3%
    % \end{theoremEnd}%
  % }{}
  \NewDocumentEnvironment{#2}{O{}O{}}{%
    \begin{theoremEnd}[#1,##2]{#3}[##1]%
    }{\end{theoremEnd}}%
  \NewDocumentEnvironment{#2D}{O{}O{}}{%
    \bgroup%
    \pratendEnableDebugSynctex{}%
    \begin{theoremEnd}[#1,##2]{#3}[##1]%
    }{\end{theoremEnd}%
    \egroup%
  }%
}

% \NewDocumentCommand{\newEndThm}{O{}mm}{%
%   % \expandafter\NewDocumentEnvironment{#2}{O{}O{}+b}{%
%     % \begin{theoremEnd}[#1,##2]{#3}[##1]%
%       % ##3%
%     % \end{theoremEnd}%
%   % }{}
%   \expandafter\NewDocumentEnvironment{#2}{O{}O{}+b}{%
%     \begin{theoremEnd}[#1,##2]{#3}[##1]%
%       ##3
%     \end{theoremEnd}%
%   }{}
% }

%%   Avoid to manually type \NewDocumentEnvironment...
%%   Usage: \newEndProof[default options]{new proof environment}{old proof environment}
%   E.g.: \newEndProof[]{proofE}{proof}
\NewDocumentCommand{\newEndProof}{O{}m}{%
  \NewDocumentEnvironment{#2}{O{}+b}{%
    \begin{proofEnd}[#1,##1]%
      ##2%
    \end{proofEnd}%
  }{}%
  \NewDocumentEnvironment{#2D}{O{}}{%
    \begin{proofEndDebug}[##1]%
    }{%
    \end{proofEndDebug}%
  }%
}

%%%%%%%%%%%%%%%%%%%%%%%%%%%%%%%%%%%%%%%%%%%%%%%%%%%%%%%%%%%%%%%%%%%%%%%
%%% Create even more automatically wrapper with option "createShortEnv"
%%%%%%%%%%%%%%%%%%%%%%%%%%%%%%%%%%%%%%%%%%%%%%%%%%%%%%%%%%%%%%%%%%%%%%%
\ifpratendOptcreateShortEnv
  %%% We create automatically the environments.
  \newEndProof{proofE}
  \newEndThm{thmE}{thm}
  \newEndThm{lemmaE}{lemma}
  \newEndThm{theoremE}{theorem}
  \newEndThm{corollaryE}{corollary}
  \newEndThm{propositionE}{proposition}
  \newEndThm{propertyE}{proposition}
  \newEndThm{factE}{proposition}
\fi

%%%%%%%%%%%%%%%%%%%%%%%%%%%%%%%%%%%%%%%%%%%%%%%%%%%%%%%%%%%%%%%%%%%%%%%
%%% Create functions to easily overwrite defaults
%%%%%%%%%%%%%%%%%%%%%%%%%%%%%%%%%%%%%%%%%%%%%%%%%%%%%%%%%%%%%%%%%%%%%%%

\NewDocumentCommand\pratendSetGlobal{m}{%
  \pgfkeys{/prAtEnd/global custom defaults/.style={
      #1
    }%
  }%
}

\NewDocumentCommand\pratendSetLocal{m}{%
  \pgfkeys{/prAtEnd/local custom defaults/.style={
      #1
    }%
  }%
}

%%%%%%%%%%%%%%%%%%%%%%%%%%%%%%%%%%%%%%%%%%%%%%%%%%%%%%%%%%%%%%%%%%%%%%%
%%% Create functions to load the category file in another file (external bibliography)
%%%%%%%%%%%%%%%%%%%%%%%%%%%%%%%%%%%%%%%%%%%%%%%%%%%%%%%%%%%%%%%%%%%%%%%

% Takes an optional category, and the job name of the external file.
\NewDocumentCommand\includeExternalAppendix{O{defaultcategory}m}{%
  \expandafter\input{\pratendGeneratePrefixFile{#2}#1}
}

\endinput
%\iffalse
%    \begin{macrocode}
%</proof-at-the-end.sty>
%    \end{macrocode}
%\fi
%\Finale
\endinput
